

% ============================================================
%  Capítulo: Resultados
% ============================================================
\chapter{Resultados}
\label{cap:resultados}

Este capítulo presenta los resultados de la inferencia de simetría sobre los 850 objetos
de cada tipo, primero para la simetría axial (Sección~\ref{sec:sm_res_eje}) y
luego para la planar (Sección~\ref{sec:sm_res_plano}). 



\newpage


\newpage
\subsubsection{Referencias sin información}
\label{sec:sm_azar}

Para interpretar un error angular hace falta saber cuánto obtendría un
predictor que no mira la imagen.

\paragraph{Eje: dirección uniforme al azar.} Para una dirección uniforme en la
esfera, el ángulo sin signo $\varphi\in[0°,90°]$ respecto de un eje fijo tiene
densidad $\sin\varphi$. De ahí:
\begin{align}
  \mathbb{E}[\varphi] &= \int_0^{\pi/2}\varphi\sin\varphi\,d\varphi = 1\ \text{rad} = 57.30°,
  \qquad \text{mediana} = 60°, \\
  \sigma &= \sqrt{\pi-3}\ \text{rad} = 21.56°, \\
  \text{P@}\theta &= 1-\cos\theta,
  \qquad \text{AUC}_{45} = 1-\frac{\sin(\pi/4)}{\pi/4} = 0.0997.
\end{align}
Así, P@5°$=0.0038$, P@10°$=0.0152$ y P@15°$=0.0341$. Como los objetos sin
predicción reciben $90°$, el valor esperado bajo azar para una fila con $N_{obj}$
objetos predichos es $(N_{obj}\cdot 57.30 + (850-N_{obj})\cdot 90)/850$, y su error
estándar $21.56\sqrt{N_{obj}}/850$. El estadístico $z$ de las tablas es la
diferencia entre el error medio observado y ese valor esperado, dividida por el
error estándar: $z<-2$ indica mejor que el azar y $z>2$, peor.

\paragraph{Plano: predictores triviales.} Para plano se usan predictores que
siempre devuelven el mismo conjunto de planos del mundo que pasan por el
origen (``siempre $X$'', ``siempre $\{X,Z\}$'', etc.), evaluados con las mismas
funciones de \texttt{evaluate.py}
(\texttt{Pipeline\_Experiments/diagnostics/trivial\_baselines.py}). Su
pertinencia se explica en la sección~\ref{sec:sm_orientacion}.

\newpage
\section{Simetría axial}
\label{sec:sm_res_eje}

\newpage
\subsection{Método azaroso}

Para saber si las predicciones aportan información sobre la simetría, se comparan con
un predictor que no mira las imágenes: la dirección del eje se elige al azar, de
manera uniforme sobre la esfera. Como el error angular se mide sin considerar el
sentido del vector ($\theta=\arccos|\hat{\mathbf{n}}\cdot\mathbf{n}^*|$), sus
valores esperados tienen forma cerrada; la Tabla~\ref{tab:axis_azar} los reporta
junto con una simulación sobre los 850 objetos y su intervalo de confianza. Un
método que no supere estos valores no extrae de los puntos de Molmo2 información
sobre el eje.



\begin{table}[htbp]
\centering
\caption{Línea base aleatoria para simetría axial: la dirección del eje se elige uniforme en la esfera, sin mirar las imágenes, para los 850 objetos (todos reciben predicción). Fórmula cerrada: valor teórico con un eje verdadero por objeto. Simulación: media de 1000 conjuntos de datos simulados, evaluados con las mismas métricas que el método; IC 95\,\%: percentiles 2.5 y 97.5.}
\label{tab:axis_azar}
\begin{adjustbox}{max width=\textwidth}
\begin{tabular}{lrrr}
\toprule
Métrica & Fórmula cerrada & Simulación & IC 95\,\% \\
\midrule
AUC$_{45^\circ}$ & 0.100 & 0.100 & [0.085, 0.113] \\
Error angular medio ($^\circ$) & 57.3 & 57.3 & [55.8, 58.7] \\
Error angular mediano ($^\circ$) & 60.0 & 60.0 & [57.7, 62.1] \\
Desv.\ estándar del error ($^\circ$) & 21.6 & 21.5 & [20.7, 22.3] \\
Precisión @5$^\circ$ & 0.004 & 0.004 & [0.000, 0.009] \\
Precisión @10$^\circ$ & 0.015 & 0.015 & [0.007, 0.024] \\
Precisión @15$^\circ$ & 0.034 & 0.034 & [0.022, 0.047] \\
\bottomrule
\end{tabular}
\end{adjustbox}
\end{table}




\newpage

\clearpage
\subsection{Flujo A: Base}

\begin{table}[!h]
\centering
\caption{Métricas base del Flujo A en simetría axial, resolución 224px, iluminación \textit{flat}, envío multiimagen.}
\label{tab:nov_axis_base}
\footnotesize
\setlength{\tabcolsep}{3pt}

\begin{tabular*}{\textwidth}{@{\extracolsep{\fill}}lrrrrrrrr@{}}
\toprule
\multirow{2}{*}{Prompt} & \multirow{2}{*}{Vistas} & \multirow{2}{*}{AUC$_{45^\circ}$} & \multicolumn{3}{c}{Error angular ($^\circ$)} & \multicolumn{3}{c}{Precisión} \\
\cmidrule(lr){4-6} \cmidrule(lr){7-9}
 &  &  & Media & Mediana & Desv. & @5$^\circ$ & @10$^\circ$ & @15$^\circ$ \\
\midrule
\multirow{3}{*}{\texttt{v00}} & 6 & 0.445 & 36.4 & 26.9 & 33.0 & 0.279 & 0.397 & 0.447 \\
 & 14 & 0.414 & 39.8 & 32.5 & 34.3 & 0.209 & 0.368 & 0.421 \\
 & 26 & 0.378 & 40.4 & 33.3 & 31.7 & 0.086 & 0.249 & 0.369 \\
\midrule
\multirow{3}{*}{\texttt{v01}} & 6 & 0.322 & 49.0 & 61.6 & 34.8 & 0.203 & 0.265 & 0.306 \\
 & 14 & 0.256 & 52.8 & 63.4 & 32.6 & 0.158 & 0.212 & 0.236 \\
 & 26 & 0.210 & 52.6 & 58.6 & 28.6 & 0.067 & 0.131 & 0.180 \\
\midrule
\multirow{3}{*}{\texttt{v02}} & 6 & 0.377 & 46.5 & 57.3 & 37.0 & 0.266 & 0.337 & 0.380 \\
 & 14 & 0.319 & 51.1 & 65.0 & 36.3 & 0.228 & 0.301 & 0.324 \\
 & 26 & 0.276 & 50.4 & 59.2 & 32.5 & 0.117 & 0.219 & 0.269 \\
\midrule
\multirow{3}{*}{\texttt{v03}} & 6 & 0.328 & 50.5 & 68.0 & 36.4 & 0.249 & 0.294 & 0.325 \\
 & 14 & 0.283 & 51.6 & 63.2 & 33.9 & 0.188 & 0.246 & 0.273 \\
 & 26 & 0.224 & 54.0 & 62.8 & 30.7 & 0.072 & 0.146 & 0.193 \\
\midrule
\multirow{3}{*}{\texttt{v04}} & 6 & \textbf{0.590} & \textbf{25.8} & \textbf{7.4} & 30.4 & \textbf{0.412} & \textbf{0.554} & 0.599 \\
 & 14 & 0.581 & 26.4 & 8.1 & 30.5 & 0.324 & 0.549 & \textbf{0.608} \\
 & 26 & 0.506 & 30.9 & 14.6 & 30.2 & 0.131 & 0.358 & 0.509 \\
\midrule
\multirow{3}{*}{\texttt{v05}} & 6 & 0.356 & 40.4 & 37.9 & 30.3 & 0.140 & 0.258 & 0.327 \\
 & 14 & 0.312 & 43.9 & 45.0 & 30.4 & 0.121 & 0.229 & 0.283 \\
 & 26 & 0.222 & 49.3 & 51.9 & 27.3 & 0.048 & 0.123 & 0.171 \\
\bottomrule
\end{tabular*}

\vspace{2ex}

\begin{tabular*}{\textwidth}{@{\extracolsep{\fill}}lrrrrrrrrrr@{}}
\toprule
\multirow{2}{*}{Prompt} & \multirow{2}{*}{Vistas} & \multicolumn{2}{c}{Err.\ trasl.} & \multicolumn{2}{c}{Err.\ trasl.\ norm.} & \multicolumn{3}{c}{SDE$_{\mathrm{ref}}$} & \multicolumn{2}{c}{Cobertura} \\
\cmidrule(lr){3-4} \cmidrule(lr){5-6} \cmidrule(lr){7-9} \cmidrule(lr){10-11}
 &  & Media & Mediana & Media & Mediana & Media & Mín. & Máx. & Objetos & Puntos \\
\midrule
\multirow{3}{*}{\texttt{v00}} & 6 & 0.4862 & 0.0513 & 0.4509 & 0.0490 & 0.0218 & 0.0000 & 2.4909 & 793 & 5.9 \\
 & 14 & 0.3005 & 0.0441 & 0.2789 & 0.0407 & 0.0161 & 0.0000 & 0.5643 & 745 & 13.6 \\
 & 26 & 0.2649 & 0.0681 & 0.2468 & 0.0607 & 0.0206 & 0.0001 & 0.4936 & 781 & 25.8 \\
\midrule
\multirow{3}{*}{\texttt{v01}} & 6 & 0.1065 & 0.0522 & 0.0970 & 0.0495 & 0.0271 & 0.0000 & 0.7157 & 818 & 5.9 \\
 & 14 & 0.0699 & 0.0426 & 0.0638 & 0.0395 & 0.0393 & 0.0000 & 0.6714 & 811 & 13.6 \\
 & 26 & 0.0902 & 0.0637 & 0.0827 & 0.0592 & 0.0582 & 0.0001 & 2.8229 & 823 & 25.4 \\
\midrule
\multirow{3}{*}{\texttt{v02}} & 6 & 0.1026 & 0.0484 & 0.0943 & 0.0458 & 0.0336 & 0.0000 & 1.0273 & 816 & 5.9 \\
 & 14 & 0.0655 & 0.0350 & 0.0602 & 0.0321 & 0.0384 & 0.0000 & 0.5774 & 773 & 13.8 \\
 & 26 & 0.0942 & 0.0658 & 0.0864 & 0.0605 & 0.0688 & 0.0001 & 2.8116 & 782 & 25.6 \\
\midrule
\multirow{3}{*}{\texttt{v03}} & 6 & 0.0974 & 0.0424 & 0.0893 & 0.0394 & 0.0366 & 0.0000 & 0.7590 & 805 & 5.8 \\
 & 14 & \textbf{0.0605} & 0.0325 & \textbf{0.0555} & 0.0293 & 0.0422 & 0.0000 & 0.7862 & 798 & 13.5 \\
 & 26 & 0.0829 & 0.0543 & 0.0762 & 0.0503 & 0.0686 & 0.0000 & 2.8219 & 764 & 25.5 \\
\midrule
\multirow{3}{*}{\texttt{v04}} & 6 & 0.2791 & 0.0269 & 0.2636 & 0.0249 & 0.0123 & 0.0000 & 0.3287 & \textbf{831} & 6.0 \\
 & 14 & 0.1452 & \textbf{0.0230} & 0.1384 & \textbf{0.0210} & \textbf{0.0101} & 0.0000 & 0.2873 & 816 & 13.8 \\
 & 26 & 0.1883 & 0.0455 & 0.1770 & 0.0418 & 0.0217 & 0.0001 & 1.5234 & 806 & 25.8 \\
\midrule
\multirow{3}{*}{\texttt{v05}} & 6 & 0.6371 & 0.1076 & 0.5969 & 0.0968 & 0.0156 & 0.0000 & 0.2020 & 818 & 6.0 \\
 & 14 & 0.4220 & 0.1176 & 0.3941 & 0.1100 & 0.0175 & 0.0001 & 0.2414 & 801 & 13.8 \\
 & 26 & 0.4465 & 0.1680 & 0.4141 & 0.1505 & 0.0232 & 0.0001 & 0.3093 & 830 & 25.8 \\
\bottomrule
\end{tabular*}

\vspace{1.5ex}
\begin{minipage}{\textwidth}
Nota: Imputación y cobertura según la Sección~\ref{sec:metricas}.
\end{minipage}
\end{table}





\newpage
\clearpage
\subsection{Flujo A: Variantes de los prompts base}

\begin{table}[H]
\centering
\caption{Métricas de las variantes \texttt{\_1} del Flujo A en simetría axial, resolución 224px, iluminación \textit{flat}, envío multiimagen.}
\label{tab:nov_axis_variantes_1}
\footnotesize
\setlength{\tabcolsep}{3pt}
\renewcommand{\arraystretch}{0.95}

\begin{tabular*}{\textwidth}{@{\extracolsep{\fill}}lrrrrrrrr@{}}
\toprule
\multirow{2}{*}{Prompt} & \multirow{2}{*}{Vistas} & \multirow{2}{*}{AUC$_{45^\circ}$} & \multicolumn{3}{c}{Error angular ($^\circ$)} & \multicolumn{3}{c}{Precisión} \\
\cmidrule(lr){4-6} \cmidrule(lr){7-9}
 &  &  & Media & Mediana & Desv. & @5$^\circ$ & @10$^\circ$ & @15$^\circ$ \\
\midrule
\multirow{3}{*}{\texttt{v00\_1}} & 6 & 0.201 & 52.1 & 56.6 & 28.0 & 0.066 & 0.111 & 0.159 \\
 & 14 & 0.185 & 53.1 & 56.8 & 27.4 & 0.059 & 0.105 & 0.137 \\
 & 26 & 0.167 & 53.2 & 56.6 & 25.5 & 0.031 & 0.072 & 0.108 \\
\midrule
\multirow{3}{*}{\texttt{v01\_1}} & 6 & 0.335 & 47.6 & 57.0 & 34.8 & 0.219 & 0.287 & 0.320 \\
 & 14 & 0.285 & 49.7 & 57.5 & 32.9 & 0.169 & 0.238 & 0.267 \\
 & 26 & 0.222 & 51.6 & 58.6 & 28.9 & 0.086 & 0.147 & 0.185 \\
\midrule
\multirow{3}{*}{\texttt{v02\_1}} & 6 & 0.429 & 40.5 & 32.2 & 36.2 & 0.281 & 0.392 & 0.435 \\
 & 14 & 0.416 & 41.5 & 34.3 & 36.1 & 0.260 & 0.371 & 0.415 \\
 & 26 & 0.297 & 51.5 & 56.7 & 35.2 & 0.097 & 0.208 & 0.271 \\
\midrule
\multirow{3}{*}{\texttt{v03\_1}} & 6 & 0.344 & 48.7 & 60.4 & 36.3 & 0.253 & 0.300 & 0.331 \\
 & 14 & 0.284 & 51.5 & 63.6 & 33.9 & 0.194 & 0.254 & 0.274 \\
 & 26 & 0.207 & 54.5 & 62.7 & 29.7 & 0.078 & 0.140 & 0.176 \\
\midrule
\multirow{3}{*}{\texttt{v04\_1}} & 6 & \textbf{0.560} & \textbf{27.9} & \textbf{8.4} & 31.2 & \textbf{0.367} & \textbf{0.523} & \textbf{0.577} \\
 & 14 & 0.545 & 29.0 & 9.2 & 31.5 & 0.319 & 0.522 & 0.567 \\
 & 26 & 0.483 & 31.6 & 15.3 & 29.5 & 0.127 & 0.358 & 0.492 \\
\midrule
\multirow{3}{*}{\texttt{v05\_1}} & 6 & 0.210 & 50.8 & 54.0 & 27.8 & 0.074 & 0.120 & 0.159 \\
 & 14 & 0.184 & 53.2 & 56.8 & 27.3 & 0.046 & 0.106 & 0.135 \\
 & 26 & 0.134 & 55.4 & 58.1 & 23.9 & 0.018 & 0.046 & 0.071 \\
\bottomrule
\end{tabular*}

\vspace{2ex}

\begin{tabular*}{\textwidth}{@{\extracolsep{\fill}}lrrrrrrrrrr@{}}
\toprule
\multirow{2}{*}{Prompt} & \multirow{2}{*}{Vistas} & \multicolumn{2}{c}{Err.\ trasl.} & \multicolumn{2}{c}{Err.\ trasl.\ norm.} & \multicolumn{3}{c}{SDE$_{\mathrm{ref}}$} & \multicolumn{2}{c}{Cobertura} \\
\cmidrule(lr){3-4} \cmidrule(lr){5-6} \cmidrule(lr){7-9} \cmidrule(lr){10-11}
 &  & Media & Mediana & Media & Mediana & Media & Mín. & Máx. & Objetos & Puntos \\
\midrule
\multirow{3}{*}{\texttt{v00\_1}} & 6 & 1.0422 & 0.3724 & 0.9629 & 0.3348 & 0.0237 & 0.0000 & 0.2951 & 791 & 5.9 \\
 & 14 & 0.6640 & 0.3245 & 0.6151 & 0.3000 & 0.0259 & 0.0000 & 0.2319 & 789 & 13.8 \\
 & 26 & 0.5968 & 0.2763 & 0.5522 & 0.2575 & 0.0296 & 0.0000 & 0.4220 & 828 & 25.8 \\
\midrule
\multirow{3}{*}{\texttt{v01\_1}} & 6 & 0.1049 & 0.0510 & 0.0965 & 0.0493 & 0.0228 & 0.0000 & 0.5640 & 827 & 5.9 \\
 & 14 & 0.0631 & 0.0361 & 0.0577 & 0.0339 & 0.0317 & 0.0000 & 0.5874 & 821 & 13.7 \\
 & 26 & 0.0935 & 0.0638 & 0.0860 & 0.0594 & 0.0703 & 0.0001 & 3.9399 & \textbf{832} & 25.5 \\
\midrule
\multirow{3}{*}{\texttt{v02\_1}} & 6 & 0.3775 & 0.0260 & 0.3495 & 0.0239 & 0.0114 & 0.0000 & 0.2990 & 678 & 5.9 \\
 & 14 & 0.1945 & \textbf{0.0232} & 0.1797 & \textbf{0.0214} & 0.0113 & 0.0000 & 0.1600 & 675 & 13.7 \\
 & 26 & 0.2105 & 0.0382 & 0.1948 & 0.0340 & 0.0248 & 0.0000 & 0.5538 & 576 & 25.5 \\
\midrule
\multirow{3}{*}{\texttt{v03\_1}} & 6 & 0.1030 & 0.0405 & 0.0935 & 0.0376 & 0.0477 & 0.0000 & 9.9825 & 814 & 5.8 \\
 & 14 & \textbf{0.0597} & 0.0314 & \textbf{0.0549} & 0.0286 & 0.0398 & 0.0000 & 0.6947 & 800 & 13.5 \\
 & 26 & 0.0874 & 0.0501 & 0.0803 & 0.0469 & 0.0895 & 0.0000 & 11.3207 & 782 & 25.0 \\
\midrule
\multirow{3}{*}{\texttt{v04\_1}} & 6 & 0.4386 & 0.0291 & 0.4161 & 0.0267 & 0.0114 & 0.0000 & 0.3038 & 824 & 5.9 \\
 & 14 & 0.2073 & 0.0258 & 0.1951 & 0.0236 & \textbf{0.0093} & 0.0000 & 0.2210 & 809 & 13.8 \\
 & 26 & 0.2269 & 0.0491 & 0.2089 & 0.0442 & 0.0155 & 0.0001 & 1.0122 & 829 & 25.7 \\
\midrule
\multirow{3}{*}{\texttt{v05\_1}} & 6 & 1.4124 & 0.3303 & 1.3303 & 0.3051 & 0.0251 & 0.0000 & 0.5648 & 804 & 5.9 \\
 & 14 & 0.6132 & 0.2528 & 0.5658 & 0.2377 & 0.0275 & 0.0002 & 0.3491 & 793 & 13.8 \\
 & 26 & 0.4594 & 0.1570 & 0.4220 & 0.1423 & 0.0298 & 0.0000 & 0.2513 & 831 & 25.9 \\
\bottomrule
\end{tabular*}

\vspace{1.5ex}
\begin{minipage}{\textwidth}
Nota: Cada prompt es una variante del prompt base con la que se busca mejorar sus resultados. Imputación y cobertura según la Sección~\ref{sec:metricas}.
\end{minipage}
\end{table}



\newpage
\clearpage
\subsection{Flujo A: Prompts adicionales}

\begin{table}[H]
\centering
\caption{Métricas de los prompts adicionales del Flujo A en simetría axial, resolución 224px, iluminación \textit{flat}, envío multiimagen.}
\label{tab:nov_axis_adicionales}
\footnotesize
\setlength{\tabcolsep}{3pt}
\renewcommand{\arraystretch}{0.95}

\begin{tabular*}{\textwidth}{@{\extracolsep{\fill}}lrrrrrrrr@{}}
\toprule
\multirow{2}{*}{Prompt} & \multirow{2}{*}{Vistas} & \multirow{2}{*}{AUC$_{45^\circ}$} & \multicolumn{3}{c}{Error angular ($^\circ$)} & \multicolumn{3}{c}{Precisión} \\
\cmidrule(lr){4-6} \cmidrule(lr){7-9}
 &  &  & Media & Mediana & Desv. & @5$^\circ$ & @10$^\circ$ & @15$^\circ$ \\
\midrule
\multirow{3}{*}{\texttt{v06}} & 6 & 0.582 & \textbf{26.5} & 7.2 & 30.8 & 0.407 & 0.547 & 0.600 \\
 & 14 & 0.563 & 27.7 & 8.5 & 31.0 & 0.331 & 0.535 & 0.592 \\
 & 26 & 0.504 & 30.4 & 14.0 & 29.6 & 0.131 & 0.380 & 0.520 \\
\midrule
\multirow{3}{*}{\texttt{v07}} & 6 & 0.284 & 55.3 & 69.6 & 36.7 & 0.167 & 0.251 & 0.279 \\
 & 14 & 0.299 & 54.7 & 71.5 & 37.4 & 0.183 & 0.266 & 0.294 \\
 & 26 & 0.195 & 61.0 & 75.1 & 32.6 & 0.065 & 0.133 & 0.165 \\
\midrule
\multirow{3}{*}{\texttt{v08}} & 6 & \textbf{0.603} & 28.0 & \textbf{4.6} & 34.8 & \textbf{0.515} & \textbf{0.594} & \textbf{0.622} \\
 & 14 & 0.554 & 31.6 & 6.7 & 35.6 & 0.431 & 0.564 & 0.582 \\
 & 26 & 0.503 & 32.8 & 13.1 & 32.7 & 0.191 & 0.431 & 0.539 \\
\bottomrule
\end{tabular*}

\vspace{2ex}

\begin{tabular*}{\textwidth}{@{\extracolsep{\fill}}lrrrrrrrrrr@{}}
\toprule
\multirow{2}{*}{Prompt} & \multirow{2}{*}{Vistas} & \multicolumn{2}{c}{Err.\ trasl.} & \multicolumn{2}{c}{Err.\ trasl.\ norm.} & \multicolumn{3}{c}{SDE$_{\mathrm{ref}}$} & \multicolumn{2}{c}{Cobertura} \\
\cmidrule(lr){3-4} \cmidrule(lr){5-6} \cmidrule(lr){7-9} \cmidrule(lr){10-11}
 &  & Media & Mediana & Media & Mediana & Media & Mín. & Máx. & Objetos & Puntos \\
\midrule
\multirow{3}{*}{\texttt{v06}} & 6 & 0.3444 & 0.0252 & 0.3250 & 0.0227 & 0.0100 & 0.0000 & 0.2014 & 824 & 5.9 \\
 & 14 & 0.1857 & 0.0252 & 0.1744 & 0.0229 & \textbf{0.0092} & 0.0000 & 0.2193 & 812 & 13.8 \\
 & 26 & 0.1669 & 0.0481 & 0.1559 & 0.0441 & 0.0337 & 0.0000 & 11.8352 & 821 & 25.6 \\
\midrule
\multirow{3}{*}{\texttt{v07}} & 6 & 0.7941 & 0.0588 & 0.7292 & 0.0532 & 0.0184 & 0.0000 & 0.1803 & 502 & 5.8 \\
 & 14 & 0.2893 & 0.0299 & 0.2695 & 0.0269 & 0.0167 & 0.0000 & 0.3465 & 493 & 13.8 \\
 & 26 & 0.2264 & 0.0379 & 0.2104 & 0.0342 & 0.0357 & 0.0000 & 0.6218 & 500 & 25.7 \\
\midrule
\multirow{3}{*}{\texttt{v08}} & 6 & 0.0611 & 0.0216 & 0.0568 & 0.0198 & 0.0154 & 0.0000 & 0.5536 & \textbf{829} & 6.0 \\
 & 14 & \textbf{0.0512} & \textbf{0.0204} & \textbf{0.0476} & \textbf{0.0187} & 0.0167 & 0.0000 & 0.6141 & 813 & 13.9 \\
 & 26 & 0.0825 & 0.0462 & 0.0764 & 0.0422 & 0.0288 & 0.0000 & 0.7664 & \textbf{829} & 25.7 \\
\bottomrule
\end{tabular*}

\vspace{1.5ex}
\begin{minipage}{\textwidth}
Nota: Imputación y cobertura según la Sección~\ref{sec:metricas}.
\end{minipage}
\end{table}





\newpage
\clearpage
\subsection{Flujo A: Top}
\label{subsec:axis_top_224}

Para elegir los mejores \textit{prompts} del Flujo A se aplica un ranking con tolerancia. Las configuraciones (\textit{prompt} por cantidad de vistas) se comparan según la métrica de mayor prioridad y se forma un grupo de empate con las que difieren del mejor valor en menos que la brecha de esa métrica. Si el grupo tiene una sola configuración, esta ocupa la posición; si no, se filtra con la métrica siguiente. Luego se repite con las configuraciones restantes. Las brechas de AUC y de precisión equivalen a unos dos errores estándar de la diferencia entre dos configuraciones evaluadas sobre $N=850$ objetos (Ecuación~\eqref{eq:se_diferencia}), mientras que las demás son heurísticas.
% Ecuacion agregada porque la referencia apuntaba a una etiqueta inexistente; revisar.
{\color{revcolor}
\begin{equation}
\label{eq:se_diferencia}
\mathrm{SE}_{\Delta} = \sqrt{\frac{2\sigma^2}{N}} \leq \sqrt{\frac{1}{2N}} \approx 0.024 ,
\end{equation}
}
\rev{donde $\sigma^2 \leq 1/4$ acota la varianza por objeto, con configuraciones independientes.}

En simetría axial, el orden de prioridad, con la brecha entre paréntesis, es: (1) \textbf{AUC} ($0.05$); (2) \textbf{error angular mediano} ($2°$); (3) \textbf{P@$5°$} ($0.05$); (4) \textbf{P@$10°$} ($0.05$); (5) \textbf{cobertura} (17 objetos, 2\% de 850); (6) \textbf{error de traslación normalizado mediano} ($0.01$), interpretable solo con la dirección correcta; y (7) \textbf{SDE$_{\mathrm{ref}}$ medio}, sin brecha, que resuelve los empates restantes.

La selección aplica este procedimiento a los promedios sobre 6, 14 y 26 vistas (Tabla~\ref{tab:axis_top_prompts}). 

\begin{table}[H]
\centering
\caption{Promedio por \textit{prompt} sobre 6, 14 y 26 vistas, en las mismas condiciones.}
\label{tab:axis_top_prompts}
\footnotesize
\setlength{\tabcolsep}{3pt}
\renewcommand{\arraystretch}{0.95}

\begin{tabular*}{\textwidth}{@{\extracolsep{\fill}}rlcrrrrrrr@{}}
\toprule
\# & Prompt & Crit. & AUC$_{45^\circ}$ & \begin{tabular}[b]{@{}c@{}}Mediana \\ ($^\circ$)\end{tabular} & P@5$^\circ$ & P@10$^\circ$ & Objetos & \begin{tabular}[b]{@{}c@{}}Trasl. \\ norm.\end{tabular} & SDE$_{\mathrm{ref}}$ \\
\midrule
1 & \texttt{v08} & 3 & 0.554 & \textbf{8.1} & \textbf{0.379} & \textbf{0.529} & \textbf{824} & \textbf{0.0269} & 0.0203 \\
2 & \texttt{v04} & 7 & \textbf{0.559} & 10.0 & 0.289 & 0.487 & 818 & 0.0292 & 0.0147 \\
3 & \texttt{v06} & 1 & 0.550 & 9.9 & 0.289 & 0.487 & 819 & 0.0299 & 0.0177 \\
4 & \texttt{v00} & 2 & 0.412 & 30.9 & 0.191 & 0.338 & 773 & 0.0501 & 0.0195 \\
-- & \texttt{v04\_1} & -- & 0.530 & 11.0 & 0.271 & 0.468 & 821 & 0.0315 & \textbf{0.0121} \\
\bottomrule
\end{tabular*}

\vspace{0.5ex}
\begin{minipage}{\textwidth}
Nota: Crit.: criterio que resolvió la posición, según el orden de las columnas. \texttt{v04\_1} 
\end{minipage}
\end{table}


Se seleccionan \texttt{v08}, \texttt{v04} y \texttt{v06}. Empatan en AUC y mediana; \texttt{v08} se separa por P@$5°$, y entre \texttt{v04} y \texttt{v06} decide solo el SDE$_{\mathrm{ref}}$, por lo que ese orden es débil. 







\newpage
\subsection{Flujo A: Experimento de resolución}

% Basado en el criterio de la sección \ref{subsec:axis_top_224} y los resultados de la tabla \ref{tab:axis_top_224}, se seleccionaron los siguientes \textit{prompts}: \texttt{v08}, \texttt{v04} y \texttt{v06}.
% Estos tres \textit{prompts} utilizan estrategias distintas, obteniendo los mayores AUC promedio (0.550,
% 0.554 y 0.547), con una clara separación del resto (el siguiente, \texttt{v00}, obtiene
% 0.408). 
% 
% Se excluye \texttt{v04\_1}, cuyo AUC promedio (0.525) también es cercano,
% por ser una variante de \texttt{v04} con rendimiento inferior en las tres
% configuraciones de vistas. Las diferencias de AUC entre los tres seleccionados no
% superan la brecha de empate, y el \textit{prompt} \texttt{v08} se distingue únicamente por su mayor precisión fina
% (P@$5°=0.502$ con $n_v=6$).


\begin{table}[H]
\centering
\caption{Métricas de los prompts seleccionados en simetría axial, resolución 448px, iluminación \textit{flat}, envío multiimagen.}
\label{tab:nov_axis_448}
\footnotesize
\setlength{\tabcolsep}{3pt}
\renewcommand{\arraystretch}{0.95}

\begin{tabular*}{\textwidth}{@{\extracolsep{\fill}}lrrrrrrrr@{}}
\toprule
\multirow{2}{*}{Prompt} & \multirow{2}{*}{Vistas} & \multirow{2}{*}{AUC$_{45^\circ}$} & \multicolumn{3}{c}{Error angular ($^\circ$)} & \multicolumn{3}{c}{Precisión} \\
\cmidrule(lr){4-6} \cmidrule(lr){7-9}
 &  &  & Media & Mediana & Desv. & @5$^\circ$ & @10$^\circ$ & @15$^\circ$ \\
\midrule
\multirow{3}{*}{\texttt{v04}} & 6 & \textbf{0.608} & \textbf{25.2} & \textbf{6.3} & 31.1 & 0.458 & \textbf{0.571} & \textbf{0.614} \\
 & 14 & 0.546 & 29.9 & 9.1 & 33.0 & 0.387 & 0.511 & 0.545 \\
 & 26 & 0.292 & 46.6 & 45.2 & 31.3 & 0.109 & 0.180 & 0.235 \\
\midrule
\multirow{3}{*}{\texttt{v06}} & 6 & 0.588 & 26.8 & 6.6 & 32.1 & 0.454 & 0.552 & 0.595 \\
 & 14 & 0.533 & 30.5 & 11.8 & 32.8 & 0.379 & 0.486 & 0.539 \\
 & 26 & 0.300 & 46.1 & 47.1 & 31.5 & 0.132 & 0.216 & 0.255 \\
\midrule
\multirow{3}{*}{\texttt{v08}} & 6 & 0.570 & 30.4 & 6.7 & 35.4 & \textbf{0.460} & 0.549 & 0.578 \\
 & 14 & 0.471 & 38.6 & 23.7 & 37.2 & 0.382 & 0.449 & 0.474 \\
 & 26 & 0.259 & 49.3 & 52.9 & 30.7 & 0.138 & 0.175 & 0.213 \\
\bottomrule
\end{tabular*}

\vspace{2ex}

\begin{tabular*}{\textwidth}{@{\extracolsep{\fill}}lrrrrrrrrrr@{}}
\toprule
\multirow{2}{*}{Prompt} & \multirow{2}{*}{Vistas} & \multicolumn{2}{c}{Err.\ trasl.} & \multicolumn{2}{c}{Err.\ trasl.\ norm.} & \multicolumn{3}{c}{SDE$_{\mathrm{ref}}$} & \multicolumn{2}{c}{Cobertura} \\
\cmidrule(lr){3-4} \cmidrule(lr){5-6} \cmidrule(lr){7-9} \cmidrule(lr){10-11}
 &  & Media & Mediana & Media & Mediana & Media & Mín. & Máx. & Objetos & Puntos \\
\midrule
\multirow{3}{*}{\texttt{v04}} & 6 & 0.1684 & 0.0249 & 0.1568 & 0.0226 & 0.0136 & 0.0000 & 0.4970 & \textbf{831} & 6.0 \\
 & 14 & 0.1576 & 0.0278 & 0.1464 & 0.0265 & 0.0141 & 0.0000 & 0.3294 & 800 & 13.8 \\
 & 26 & 0.2741 & 0.0617 & 0.2537 & 0.0575 & 0.0402 & 0.0000 & 5.6395 & 750 & 25.8 \\
\midrule
\multirow{3}{*}{\texttt{v06}} & 6 & 0.2480 & 0.0240 & 0.2331 & 0.0226 & \textbf{0.0127} & 0.0000 & 0.3046 & 818 & 5.9 \\
 & 14 & 0.1940 & 0.0271 & 0.1792 & 0.0251 & 0.0142 & 0.0000 & 0.4193 & 799 & 13.8 \\
 & 26 & 0.2682 & 0.0508 & 0.2455 & 0.0462 & 0.0321 & 0.0000 & 2.5148 & 766 & 25.7 \\
\midrule
\multirow{3}{*}{\texttt{v08}} & 6 & 0.0678 & \textbf{0.0199} & 0.0614 & \textbf{0.0185} & 0.0198 & 0.0000 & 0.4583 & 821 & 5.9 \\
 & 14 & \textbf{0.0656} & 0.0216 & \textbf{0.0603} & 0.0198 & 0.0305 & 0.0000 & 0.7506 & 771 & 13.8 \\
 & 26 & 0.1055 & 0.0551 & 0.0963 & 0.0512 & 0.0544 & 0.0001 & 2.7252 & 779 & 25.5 \\
\bottomrule
\end{tabular*}

\vspace{1.5ex}
\begin{minipage}{\textwidth}
Nota: Imputación y cobertura según la Sección~\ref{sec:sm_cobertura}.
\end{minipage}
\end{table}





A 448px, \texttt{v04} y \texttt{v06} redujeron su AUC promedio en una magnitud similar respecto de 224px, 
a diferencia de
\texttt{v08}, cuya mediana con 14 vistas también aumentó de forma marcada. Por tanto, para el  
texttt{v04} se eligió para el experimento a 1136px y posteriores. 





\newpage

\clearpage
\begin{table}[H]
\centering
\caption{Métricas del prompt seleccionado en simetría axial, resolución 1136px, iluminación \textit{flat}, envío multiimagen.}
\label{tab:nov_axis_1136}
\footnotesize
\setlength{\tabcolsep}{3pt}
\renewcommand{\arraystretch}{0.95}

\begin{tabular*}{\textwidth}{@{\extracolsep{\fill}}lrrrrrrrr@{}}
\toprule
\multirow{2}{*}{Prompt} & \multirow{2}{*}{Vistas} & \multirow{2}{*}{AUC$_{45^\circ}$} & \multicolumn{3}{c}{Error angular ($^\circ$)} & \multicolumn{3}{c}{Precisión} \\
\cmidrule(lr){4-6} \cmidrule(lr){7-9}
 &  &  & Media & Mediana & Desv. & @5$^\circ$ & @10$^\circ$ & @15$^\circ$ \\
\midrule
\multirow{3}{*}{\texttt{v04}} & 6 & \textbf{0.577} & \textbf{27.2} & \textbf{8.4} & 31.6 & \textbf{0.421} & \textbf{0.536} & \textbf{0.592} \\
 & 14 & 0.416 & 38.3 & 30.6 & 32.8 & 0.227 & 0.326 & 0.392 \\
 & 26 & 0.188 & 52.7 & 54.7 & 27.3 & 0.035 & 0.078 & 0.123 \\
\bottomrule
\end{tabular*}

\vspace{2ex}

\begin{tabular*}{\textwidth}{@{\extracolsep{\fill}}lrrrrrrrrrr@{}}
\toprule
\multirow{2}{*}{Prompt} & \multirow{2}{*}{Vistas} & \multicolumn{2}{c}{Err.\ trasl.} & \multicolumn{2}{c}{Err.\ trasl.\ norm.} & \multicolumn{3}{c}{SDE$_{\mathrm{ref}}$} & \multicolumn{2}{c}{Cobertura} \\
\cmidrule(lr){3-4} \cmidrule(lr){5-6} \cmidrule(lr){7-9} \cmidrule(lr){10-11}
 &  & Media & Mediana & Media & Mediana & Media & Mín. & Máx. & Objetos & Puntos \\
\midrule
\multirow{3}{*}{\texttt{v04}} & 6 & 0.3635 & \textbf{0.0261} & 0.3423 & \textbf{0.0242} & \textbf{0.0133} & 0.0000 & 0.3934 & \textbf{816} & 5.9 \\
 & 14 & \textbf{0.2835} & 0.0546 & \textbf{0.2613} & 0.0494 & 0.0199 & 0.0000 & 0.3007 & 775 & 13.9 \\
 & 26 & 0.3007 & 0.0952 & 0.2780 & 0.0878 & 0.0345 & 0.0001 & 1.0829 & 765 & 25.9 \\
\bottomrule
\end{tabular*}

\vspace{1.5ex}
\begin{minipage}{\textwidth}
Nota: Imputación y cobertura según la Sección~\ref{sec:sm_cobertura}.
\end{minipage}
\end{table}



\newpage
\subsection{Flujo A: Experimento de iluminación}

% Este experimento evalúa si el desempeño depende de la exposición de los renders. Se
% mantienen fijos el \textit{prompt} (\texttt{v04}), la resolución (224px) y las
% vistas, y se cambia solo el tono de gris uniforme de la malla: \textit{darker}
% (0.3), \textit{flat} (0.7, la condición usada en el resto del capítulo) y
% \textit{brighter} (0.95). Las tres condiciones varían la intensidad, no la dirección
% de la luz, por lo que no evalúan si el sombreado ayuda a percibir la geometría, sino
% si el método es robusto a la exposición. Se eligió 224px porque es la única
% resolución en que \texttt{v04} se mantiene estable con 6, 14 y 26 vistas, lo que
% evita mezclar el efecto de la iluminación con el de la resolución.
% 
% Nota: Ojo con la redacción y las citas. En \texttt{ImagesGenerator/export\_fibonacci\_views.py}, la ``iluminación'' se implementa cambiando el color gris de la malla (TexturesVertex), con la luz direccional fija. No cambia la intensidad de la luz. Por eso escribí ``tono de gris uniforme de la malla'' y no ``intensidad de la luz''. \texttt{lectura\_iluminacion.md} advierte que 3D Common Corruptions solo aplica si se modifica la luz de la escena; con esta implementación no correspondería citarlo. VLM-RobustBench, que evalúa Molmo ante perturbaciones fotométricas, sí sirve, pero tiene la venue pendiente de verificar.


\begin{table}[H]
\centering
\caption{Efecto de la iluminación en simetría axial con el \textit{prompt} \texttt{v04}, resolución 224px, envío multiimagen.}
\label{tab:nov_axis_iluminacion}
\footnotesize
\setlength{\tabcolsep}{3pt}
\renewcommand{\arraystretch}{0.95}

\begin{tabular*}{\textwidth}{@{\extracolsep{\fill}}lrrrrrrrr@{}}
\toprule
\multirow{2}{*}{Iluminación} & \multirow{2}{*}{Vistas} & \multirow{2}{*}{AUC$_{45^\circ}$} & \multicolumn{3}{c}{Error angular ($^\circ$)} & \multicolumn{3}{c}{Precisión} \\
\cmidrule(lr){4-6} \cmidrule(lr){7-9}
 &  &  & Media & Mediana & Desv. & @5$^\circ$ & @10$^\circ$ & @15$^\circ$ \\
\midrule
\multirow{3}{*}{\textit{flat}} & 6 & 0.590 & 25.8 & \textbf{7.4} & 30.4 & 0.412 & 0.554 & 0.599 \\
 & 14 & 0.581 & 26.4 & 8.1 & 30.5 & 0.324 & 0.549 & 0.608 \\
 & 26 & 0.506 & 30.9 & 14.6 & 30.2 & 0.131 & 0.358 & 0.509 \\
\midrule
\multirow{3}{*}{\textit{darker}} & 6 & 0.589 & 26.2 & 7.5 & 30.7 & \textbf{0.413} & 0.545 & 0.602 \\
 & 14 & 0.591 & 25.9 & 7.7 & 30.4 & 0.324 & \textbf{0.573} & \textbf{0.621} \\
 & 26 & 0.508 & 30.3 & 14.8 & 29.6 & 0.119 & 0.346 & 0.511 \\
\midrule
\multirow{3}{*}{\textit{brighter}} & 6 & \textbf{0.594} & \textbf{25.4} & \textbf{7.4} & 30.2 & 0.400 & 0.556 & 0.611 \\
 & 14 & 0.581 & 26.1 & 8.7 & 29.9 & 0.326 & 0.548 & 0.609 \\
 & 26 & 0.493 & 31.4 & 15.7 & 29.8 & 0.103 & 0.334 & 0.489 \\
\bottomrule
\end{tabular*}

\vspace{2ex}

\begin{tabular*}{\textwidth}{@{\extracolsep{\fill}}lrrrrrrrrrr@{}}
\toprule
\multirow{2}{*}{Iluminación} & \multirow{2}{*}{Vistas} & \multicolumn{2}{c}{Err.\ trasl.} & \multicolumn{2}{c}{Err.\ trasl.\ norm.} & \multicolumn{3}{c}{SDE$_{\mathrm{ref}}$} & \multicolumn{2}{c}{Cobertura} \\
\cmidrule(lr){3-4} \cmidrule(lr){5-6} \cmidrule(lr){7-9} \cmidrule(lr){10-11}
 &  & Media & Mediana & Media & Mediana & Media & Mín. & Máx. & Objetos & Puntos \\
\midrule
\multirow{3}{*}{\textit{flat}} & 6 & 0.2791 & 0.0269 & 0.2636 & 0.0249 & 0.0123 & 0.0000 & 0.3287 & 831 & 6.0 \\
 & 14 & 0.1452 & \textbf{0.0230} & 0.1384 & \textbf{0.0210} & \textbf{0.0101} & 0.0000 & 0.2873 & 816 & 13.8 \\
 & 26 & 0.1883 & 0.0455 & 0.1770 & 0.0418 & 0.0217 & 0.0001 & 1.5234 & 806 & 25.8 \\
\midrule
\multirow{3}{*}{\textit{darker}} & 6 & 0.1847 & 0.0276 & 0.1730 & 0.0258 & 0.0124 & 0.0000 & 0.2400 & 831 & 5.9 \\
 & 14 & \textbf{0.1143} & 0.0250 & \textbf{0.1073} & 0.0238 & 0.0102 & 0.0000 & 0.2991 & 822 & 13.8 \\
 & 26 & 0.1405 & 0.0483 & 0.1306 & 0.0452 & 0.0209 & 0.0001 & 0.8342 & 812 & 25.7 \\
\midrule
\multirow{3}{*}{\textit{brighter}} & 6 & 0.2231 & 0.0251 & 0.2108 & 0.0225 & 0.0124 & 0.0000 & 0.5158 & \textbf{835} & 5.9 \\
 & 14 & 0.1530 & 0.0240 & 0.1451 & 0.0221 & 0.0104 & 0.0000 & 0.2866 & 824 & 13.8 \\
 & 26 & 0.1571 & 0.0494 & 0.1470 & 0.0451 & 0.0171 & 0.0000 & 0.5667 & 807 & 25.8 \\
\bottomrule
\end{tabular*}

\vspace{1.5ex}
\begin{minipage}{\textwidth}
Nota: Iluminación: valor de gris uniforme de la malla (\textit{flat} 0.7, \textit{darker} 0.3, \textit{brighter} 0.95). Imputación y cobertura según la Sección~\ref{sec:sm_cobertura}.
\end{minipage}
\end{table}

\clearpage


\newpage    
\subsection{Flujo B y C}

% Los Flujos B y C prueban si dar contexto sobre el objeto mejora el \textit{pointing}.
% Se usa \texttt{v04} a 224px con iluminación \textit{flat}. En el Flujo B, una vista
% del objeto (elegida al azar con semilla fija, con elevación entre $-60^\circ$ y
% $+60^\circ$) se describe primero con Molmo2, y esa descripción se antepone al
% mismo \textit{prompt} \texttt{v04}. En el Flujo C, esa llamada previa además
% describe en texto por dónde pasa el eje, y el \textit{prompt} principal se
% reemplaza por uno genérico que pide exactamente dos puntos sobre el eje en cada
% vista. Como ese \textit{prompt} no depende del \textit{prompt} base, el Flujo C se
% evaluó con la corrida existente, que usa \texttt{v05\_1} solo como respaldo cuando
% no se obtiene la pista de ubicación.


\begin{table}[H]
\centering
\caption{Comparación de los Flujos A, B y C en simetría axial con el \textit{prompt} \texttt{v04}, resolución 224px, iluminación \textit{flat}, envío multiimagen.}
\label{tab:nov_axis_flujos}
\footnotesize
\setlength{\tabcolsep}{3pt}
\renewcommand{\arraystretch}{0.95}

\begin{tabular*}{\textwidth}{@{\extracolsep{\fill}}lrrrrrrrr@{}}
\toprule
\multirow{2}{*}{Flujo} & \multirow{2}{*}{Vistas} & \multirow{2}{*}{AUC$_{45^\circ}$} & \multicolumn{3}{c}{Error angular ($^\circ$)} & \multicolumn{3}{c}{Precisión} \\
\cmidrule(lr){4-6} \cmidrule(lr){7-9}
 &  &  & Media & Mediana & Desv. & @5$^\circ$ & @10$^\circ$ & @15$^\circ$ \\
\midrule
\multirow{3}{*}{A} & 6 & 0.590 & 25.8 & 7.4 & 30.4 & 0.412 & 0.554 & 0.599 \\
 & 14 & 0.581 & 26.4 & 8.1 & 30.5 & 0.324 & 0.549 & 0.608 \\
 & 26 & 0.506 & 30.9 & 14.6 & 30.2 & 0.131 & 0.358 & 0.509 \\
\midrule
\multirow{3}{*}{B} & 6 & \textbf{0.606} & \textbf{24.7} & \textbf{7.1} & 30.0 & \textbf{0.418} & \textbf{0.575} & \textbf{0.626} \\
 & 14 & 0.575 & 26.8 & 8.5 & 30.5 & 0.311 & 0.538 & 0.601 \\
 & 26 & 0.491 & 31.2 & 15.4 & 29.5 & 0.098 & 0.339 & 0.492 \\
\midrule
\multirow{3}{*}{C} & 6 & 0.317 & 58.2 & 90.0 & 40.3 & 0.261 & 0.319 & 0.334 \\
 & 14 & 0.215 & 68.8 & 90.0 & 36.4 & 0.167 & 0.226 & 0.232 \\
 & 26 & 0.126 & 76.5 & 90.0 & 29.7 & 0.053 & 0.106 & 0.133 \\
\bottomrule
\end{tabular*}

\vspace{2ex}

\begin{tabular*}{\textwidth}{@{\extracolsep{\fill}}lrrrrrrrrrr@{}}
\toprule
\multirow{2}{*}{Flujo} & \multirow{2}{*}{Vistas} & \multicolumn{2}{c}{Err.\ trasl.} & \multicolumn{2}{c}{Err.\ trasl.\ norm.} & \multicolumn{3}{c}{SDE$_{\mathrm{ref}}$} & \multicolumn{2}{c}{Cobertura} \\
\cmidrule(lr){3-4} \cmidrule(lr){5-6} \cmidrule(lr){7-9} \cmidrule(lr){10-11}
 &  & Media & Mediana & Media & Mediana & Media & Mín. & Máx. & Objetos & Puntos \\
\midrule
\multirow{3}{*}{A} & 6 & 0.2791 & 0.0269 & 0.2636 & 0.0249 & 0.0123 & 0.0000 & 0.3287 & 831 & 6.0 \\
 & 14 & 0.1452 & 0.0230 & 0.1384 & 0.0210 & \textbf{0.0101} & 0.0000 & 0.2873 & 816 & 13.8 \\
 & 26 & 0.1883 & 0.0455 & 0.1770 & 0.0418 & 0.0217 & 0.0001 & 1.5234 & 806 & 25.8 \\
\midrule
\multirow{3}{*}{B} & 6 & 0.2361 & 0.0274 & 0.2252 & 0.0248 & 0.0111 & 0.0000 & 0.3090 & \textbf{836} & 5.9 \\
 & 14 & 0.1706 & 0.0254 & 0.1639 & 0.0231 & 0.0103 & 0.0000 & 0.2009 & 827 & 13.8 \\
 & 26 & 0.1678 & 0.0491 & 0.1563 & 0.0460 & 0.0189 & 0.0001 & 1.0167 & 823 & 25.6 \\
\midrule
\multirow{3}{*}{C} & 6 & 0.1010 & 0.0190 & 0.0930 & 0.0169 & 0.0310 & 0.0000 & 3.0304 & 374 & 4.6 \\
 & 14 & \textbf{0.0520} & \textbf{0.0148} & \textbf{0.0470} & \textbf{0.0140} & 0.0112 & 0.0000 & 0.2677 & 244 & 12.2 \\
 & 26 & 0.0894 & 0.0375 & 0.0792 & 0.0339 & 0.0421 & 0.0001 & 2.0649 & 168 & 24.6 \\
\bottomrule
\end{tabular*}

\vspace{1.5ex}
\begin{minipage}{\textwidth}
Nota: Flujo A: \textit{prompt} \texttt{v04} sin contexto. Flujo B: el mismo \textit{prompt} precedido por una descripción del objeto. Flujo C: \textit{prompt} genérico construido a partir de la descripción y de una pista de ubicación del eje; no depende del \textit{prompt} base, que solo se usa como respaldo (corrida con \texttt{v05\_1} como respaldo). Imputación y cobertura según la Sección~\ref{sec:sm_cobertura}.
\end{minipage}
\end{table}




\newpage

\subsection{Flujo A: Envío individual vs Multiimagen}

% Se repite la comparación con \texttt{v05\_1}. En plano debe considerarse además que
% el estimador se degrada con muchas vistas incluso con puntos perfectos
% (Sección~\ref{sec:puntuacion_plano}), por lo que el resultado con 26 vistas no
% puede atribuirse solo a Molmo2.
% 

\begin{table}[H]
\centering
\caption{Envío multiimagen frente a individual en simetría axial con el \textit{prompt} \texttt{v04}, resolución 224px, iluminación \textit{flat}.}
\label{tab:nov_axis_modo_envio}
\footnotesize
\setlength{\tabcolsep}{3pt}
\renewcommand{\arraystretch}{0.95}

\begin{tabular*}{\textwidth}{@{\extracolsep{\fill}}lrrrrrrrr@{}}
\toprule
\multirow{2}{*}{Envío} & \multirow{2}{*}{Vistas} & \multirow{2}{*}{AUC$_{45^\circ}$} & \multicolumn{3}{c}{Error angular ($^\circ$)} & \multicolumn{3}{c}{Precisión} \\
\cmidrule(lr){4-6} \cmidrule(lr){7-9}
 &  &  & Media & Mediana & Desv. & @5$^\circ$ & @10$^\circ$ & @15$^\circ$ \\
\midrule
\multirow{3}{*}{\textit{multi}} & 6 & \textbf{0.590} & \textbf{25.8} & \textbf{7.4} & 30.4 & \textbf{0.412} & \textbf{0.554} & 0.599 \\
 & 14 & 0.581 & 26.4 & 8.1 & 30.5 & 0.324 & 0.549 & \textbf{0.608} \\
 & 26 & 0.506 & 30.9 & 14.6 & 30.2 & 0.131 & 0.358 & 0.509 \\
\midrule
\multirow{3}{*}{\textit{single}} & 6 & 0.496 & 31.6 & 14.6 & 31.1 & 0.298 & 0.432 & 0.503 \\
 & 14 & 0.520 & 30.1 & 12.5 & 30.9 & 0.335 & 0.471 & 0.526 \\
 & 26 & 0.535 & 29.5 & 10.6 & 31.4 & 0.366 & 0.493 & 0.547 \\
\bottomrule
\end{tabular*}

\vspace{2ex}

\begin{tabular*}{\textwidth}{@{\extracolsep{\fill}}lrrrrrrrrrr@{}}
\toprule
\multirow{2}{*}{Envío} & \multirow{2}{*}{Vistas} & \multicolumn{2}{c}{Err.\ trasl.} & \multicolumn{2}{c}{Err.\ trasl.\ norm.} & \multicolumn{3}{c}{SDE$_{\mathrm{ref}}$} & \multicolumn{2}{c}{Cobertura} \\
\cmidrule(lr){3-4} \cmidrule(lr){5-6} \cmidrule(lr){7-9} \cmidrule(lr){10-11}
 &  & Media & Mediana & Media & Mediana & Media & Mín. & Máx. & Objetos & Puntos \\
\midrule
\multirow{3}{*}{\textit{multi}} & 6 & 0.2791 & 0.0269 & 0.2636 & 0.0249 & 0.0123 & 0.0000 & 0.3287 & 831 & 6.0 \\
 & 14 & 0.1452 & 0.0230 & 0.1384 & 0.0210 & 0.0101 & 0.0000 & 0.2873 & 816 & 13.8 \\
 & 26 & 0.1883 & 0.0455 & 0.1770 & 0.0418 & 0.0217 & 0.0001 & 1.5234 & 806 & 25.8 \\
\midrule
\multirow{3}{*}{\textit{single}} & 6 & 0.2229 & 0.0351 & 0.2115 & 0.0331 & 0.0161 & 0.0000 & 0.3202 & \textbf{850} & 6.0 \\
 & 14 & 0.0986 & 0.0266 & 0.0928 & 0.0248 & 0.0117 & 0.0000 & 0.2270 & \textbf{850} & 14.0 \\
 & 26 & \textbf{0.0563} & \textbf{0.0221} & \textbf{0.0526} & \textbf{0.0199} & \textbf{0.0094} & 0.0000 & 0.2213 & \textbf{850} & 25.9 \\
\bottomrule
\end{tabular*}

\vspace{1.5ex}
\begin{minipage}{\textwidth}
Nota: \textit{multi}: las $n_v$ vistas se envían en una sola llamada, con el \textit{prompt} multivista. \textit{single}: cada vista se envía por separado, con el \textit{prompt} de imagen única, y los puntos se reúnen antes de la triangulación. Imputación y cobertura según la Sección~\ref{sec:sm_cobertura}.
\end{minipage}
\end{table}







\newpage
\subsection{Comparación de modelos}

\begin{table}[H]
\centering
\caption{Comparación de VLMs en simetría axial con el \textit{prompt} \texttt{v04}, resolución 224px, iluminación \textit{flat}, envío individual.}
\label{tab:nov_axis_vlms}
\footnotesize
\setlength{\tabcolsep}{3pt}
\renewcommand{\arraystretch}{0.95}

\begin{tabular*}{\textwidth}{@{\extracolsep{\fill}}lrrrrrrrr@{}}
\toprule
\multirow{2}{*}{Modelo} & \multirow{2}{*}{Vistas} & \multirow{2}{*}{AUC$_{45^\circ}$} & \multicolumn{3}{c}{Error angular ($^\circ$)} & \multicolumn{3}{c}{Precisión} \\
\cmidrule(lr){4-6} \cmidrule(lr){7-9}
 &  &  & Media & Mediana & Desv. & @5$^\circ$ & @10$^\circ$ & @15$^\circ$ \\
\midrule
\multirow{3}{*}{Molmo2-8B} & 6 & 0.496 & 31.6 & 14.6 & 31.1 & 0.298 & 0.432 & 0.503 \\
 & 14 & 0.520 & 30.1 & 12.5 & 30.9 & 0.335 & 0.471 & 0.526 \\
 & 26 & 0.535 & 29.5 & 10.6 & 31.4 & 0.366 & 0.493 & 0.547 \\
\midrule
\multirow{3}{*}{MolmoPoint-8B} & 6 & 0.574 & 26.5 & 9.2 & 30.3 & 0.386 & 0.509 & 0.573 \\
 & 14 & 0.620 & 23.9 & 6.0 & 30.0 & 0.458 & 0.574 & 0.635 \\
 & 26 & \textbf{0.672} & \textbf{20.8} & \textbf{4.2} & 29.1 & \textbf{0.541} & \textbf{0.642} & \textbf{0.694} \\
\midrule
\multirow{3}{*}{Qwen3-VL-8B} & 6 & 0.322 & 40.8 & 35.4 & 27.3 & 0.047 & 0.132 & 0.236 \\
 & 14 & 0.331 & 40.1 & 33.5 & 27.3 & 0.048 & 0.139 & 0.253 \\
 & 26 & 0.346 & 39.8 & 33.6 & 28.3 & 0.087 & 0.154 & 0.269 \\
\bottomrule
\end{tabular*}

\vspace{2ex}

\begin{tabular*}{\textwidth}{@{\extracolsep{\fill}}lrrrrrrrrrr@{}}
\toprule
\multirow{2}{*}{Modelo} & \multirow{2}{*}{Vistas} & \multicolumn{2}{c}{Err.\ trasl.} & \multicolumn{2}{c}{Err.\ trasl.\ norm.} & \multicolumn{3}{c}{SDE$_{\mathrm{ref}}$} & \multicolumn{2}{c}{Cobertura} \\
\cmidrule(lr){3-4} \cmidrule(lr){5-6} \cmidrule(lr){7-9} \cmidrule(lr){10-11}
 &  & Media & Mediana & Media & Mediana & Media & Mín. & Máx. & Objetos & Puntos \\
\midrule
\multirow{3}{*}{Molmo2-8B} & 6 & 0.2229 & 0.0351 & 0.2115 & 0.0331 & 0.0161 & 0.0000 & 0.3202 & \textbf{850} & 6.0 \\
 & 14 & 0.0986 & 0.0266 & 0.0928 & 0.0248 & 0.0117 & 0.0000 & 0.2270 & \textbf{850} & 14.0 \\
 & 26 & 0.0563 & 0.0221 & 0.0526 & 0.0199 & 0.0094 & 0.0000 & 0.2213 & \textbf{850} & 25.9 \\
\midrule
\multirow{3}{*}{MolmoPoint-8B} & 6 & 0.0966 & 0.0212 & 0.0910 & 0.0189 & 0.0245 & 0.0000 & 1.2469 & 842 & 5.0 \\
 & 14 & 0.0390 & 0.0147 & 0.0363 & 0.0134 & 0.0169 & 0.0000 & 1.7116 & 849 & 11.6 \\
 & 26 & \textbf{0.0270} & \textbf{0.0114} & \textbf{0.0250} & \textbf{0.0106} & \textbf{0.0070} & 0.0000 & 0.3969 & \textbf{850} & 21.5 \\
\midrule
\multirow{3}{*}{Qwen3-VL-8B} & 6 & 0.5517 & 0.1775 & 0.5213 & 0.1629 & 0.0395 & 0.0003 & 0.6324 & 849 & 6.0 \\
 & 14 & 0.2888 & 0.1263 & 0.2726 & 0.1158 & 0.0356 & 0.0002 & 0.4402 & \textbf{850} & 13.9 \\
 & 26 & 0.1868 & 0.1093 & 0.1774 & 0.1001 & 0.0275 & 0.0002 & 0.5168 & \textbf{850} & 25.8 \\
\bottomrule
\end{tabular*}

\vspace{1.5ex}
\begin{minipage}{\textwidth}
Nota: Cada vista se envía en una llamada aparte y los puntos se reúnen antes de la triangulación. MolmoPoint hace dos llamadas por vista, una por punto, con una versión abreviada del \textit{prompt}; Qwen3-VL usa el mismo \textit{prompt}, adaptado a su formato de salida JSON. Imputación y cobertura según la Sección~\ref{sec:sm_cobertura}.
\end{minipage}
\end{table}


\subsection{Retroproyección}






\newpage
% ------------------------------------------------------------
\section{Simetría planar}
\label{sec:sm_res_plano}

\subsection{Método azaroso}

Se aplica la misma línea base que en la simetría axial, ahora sobre la normal del
plano, que se elige uniforme en la esfera; el plano pasa por el centro del objeto.
La Tabla~\ref{tab:plane_azar} incluye además $F1_{\mathrm{ref}}$. Como los objetos
planares conservan la orientación canónica de ShapeNet, el azar uniforme es una cota
inferior débil: la comparación exigente es contra los predictores triviales de la
Sección~\ref{sec:sm_triviales}.


\begin{table}[htbp]
\centering
\caption{Línea base aleatoria para simetría planar: la normal del plano se elige uniforme en la esfera y el plano pasa por el centro del objeto ($d=0$), sin mirar las imágenes, para los 850 objetos (todos reciben predicción; 692 con 1 plano, 152 con 2 planos y 6 con 3 planos). El error se mide contra el plano verdadero más cercano, igual que para el método, por lo que el azar es más favorable que el valor teórico de un solo plano ($57.3^\circ$). Simulación: media de 1000 conjuntos de datos simulados; IC 95\,\%: percentiles 2.5 y 97.5.}
\label{tab:plane_azar}
\begin{adjustbox}{max width=\textwidth}
\begin{tabular}{lrr}
\toprule
Métrica & Simulación & IC 95\,\% \\
\midrule
AUC$_{45^\circ}$ & 0.118 & [0.104, 0.132] \\
Error angular medio ($^\circ$) & 54.6 & [53.1, 56.1] \\
Error angular mediano ($^\circ$) & 56.0 & [53.6, 58.3] \\
Desv.\ estándar del error ($^\circ$) & 21.8 & [21.1, 22.6] \\
Precisión @5$^\circ$ & 0.004 & [0.001, 0.009] \\
Precisión @10$^\circ$ & 0.018 & [0.009, 0.027] \\
Precisión @15$^\circ$ & 0.040 & [0.028, 0.054] \\
$F1_{\mathrm{ref}}$ voraz & 0.009 & [0.005, 0.014] \\
$F1_{\mathrm{ref}}$ húngaro & 0.010 & [0.005, 0.015] \\
\bottomrule
\end{tabular}
\end{adjustbox}
\end{table}


\newpage
\clearpage
\subsection{Flujo A: Base}

\begin{table}[H]
\centering
\caption{Métricas base del Flujo A en simetría planar, resolución 224px, iluminación \textit{flat}, envío multiimagen.}
\label{tab:nov_plane_base}
\footnotesize
\setlength{\tabcolsep}{3pt}
\renewcommand{\arraystretch}{0.95}

\begin{tabular*}{\textwidth}{@{\extracolsep{\fill}}lrrrrrrrr@{}}
\toprule
\multirow{2}{*}{Prompt} & \multirow{2}{*}{Vistas} & \multirow{2}{*}{AUC$_{45^\circ}$} & \multicolumn{3}{c}{Error angular ($^\circ$)} & \multicolumn{3}{c}{Precisión} \\
\cmidrule(lr){4-6} \cmidrule(lr){7-9}
 &  &  & Media & Mediana & Desv. & @5$^\circ$ & @10$^\circ$ & @15$^\circ$ \\
\midrule
\multirow{3}{*}{\texttt{v00}} & 6 & 0.663 & 26.2 & 4.9 & 36.7 & 0.508 & 0.707 & 0.724 \\
 & 14 & 0.567 & 32.9 & 5.6 & 37.6 & 0.472 & 0.585 & 0.621 \\
 & 26 & 0.666 & 23.1 & 4.6 & 32.5 & 0.513 & 0.615 & 0.686 \\
\midrule
\multirow{3}{*}{\texttt{v01}} & 6 & 0.155 & 70.9 & 88.7 & 31.8 & 0.140 & 0.151 & 0.153 \\
 & 14 & 0.025 & 79.6 & 84.7 & 15.3 & 0.011 & 0.018 & 0.025 \\
 & 26 & 0.019 & 87.2 & 89.6 & 12.5 & 0.018 & 0.019 & 0.020 \\
\midrule
\multirow{3}{*}{\texttt{v02}} & 6 & \textbf{0.779} & 16.6 & 3.0 & 30.9 & \textbf{0.752} & \textbf{0.829} & \textbf{0.833} \\
 & 14 & 0.628 & 29.2 & 4.8 & 37.8 & 0.529 & 0.668 & 0.697 \\
 & 26 & 0.773 & \textbf{14.0} & \textbf{2.6} & 25.1 & 0.705 & 0.754 & 0.799 \\
\midrule
\multirow{3}{*}{\texttt{v03}} & 6 & 0.171 & 68.8 & 87.6 & 32.9 & 0.151 & 0.168 & 0.176 \\
 & 14 & 0.036 & 78.6 & 84.8 & 17.6 & 0.021 & 0.034 & 0.036 \\
 & 26 & 0.020 & 85.4 & 89.5 & 13.8 & 0.015 & 0.017 & 0.019 \\
\midrule
\multirow{3}{*}{\texttt{v04}} & 6 & 0.669 & 27.5 & 3.0 & 39.1 & 0.689 & 0.706 & 0.712 \\
 & 14 & 0.630 & 29.6 & 4.7 & 38.5 & 0.588 & 0.680 & 0.697 \\
 & 26 & 0.736 & 18.6 & 2.7 & 31.0 & 0.711 & 0.720 & 0.758 \\
\midrule
\multirow{3}{*}{\texttt{v05}} & 6 & 0.232 & 62.8 & 85.4 & 35.8 & 0.186 & 0.219 & 0.233 \\
 & 14 & 0.104 & 72.6 & 85.4 & 26.4 & 0.061 & 0.094 & 0.109 \\
 & 26 & 0.102 & 77.0 & 88.5 & 26.8 & 0.065 & 0.105 & 0.111 \\
\bottomrule
\end{tabular*}

\vspace{2ex}

\begin{tabular*}{\textwidth}{@{\extracolsep{\fill}}lrrrrrrrrrr@{}}
\toprule
\multirow{2}{*}{Prompt} & \multirow{2}{*}{Vistas} & \multicolumn{2}{c}{Err.\ trasl.} & \multicolumn{2}{c}{Err.\ trasl.\ norm.} & \multicolumn{3}{c}{SDE$_{\mathrm{ref}}$} & \multicolumn{2}{c}{Cobertura} \\
\cmidrule(lr){3-4} \cmidrule(lr){5-6} \cmidrule(lr){7-9} \cmidrule(lr){10-11}
 &  & Media & Mediana & Media & Mediana & Media & Mín. & Máx. & Objetos & Puntos \\
\midrule
\multirow{3}{*}{\texttt{v00}} & 6 & 0.0590 & 0.0115 & 0.0591 & 0.0115 & 0.0179 & 0.0000 & 5.6727 & 833 & 4.1 \\
 & 14 & 0.0259 & 0.0071 & 0.0260 & 0.0071 & 0.0150 & 0.0000 & 0.6917 & 777 & 7.8 \\
 & 26 & 0.0185 & 0.0022 & 0.0186 & 0.0022 & 0.0105 & 0.0000 & 0.6871 & 810 & 14.7 \\
\midrule
\multirow{3}{*}{\texttt{v01}} & 6 & 0.1348 & 0.0582 & 0.1368 & 0.0582 & 0.0172 & 0.0000 & 0.3684 & 846 & 3.1 \\
 & 14 & 0.0477 & 0.0245 & 0.0477 & 0.0249 & 0.0367 & 0.0000 & 0.3298 & 846 & 6.6 \\
 & 26 & 0.0188 & 0.0050 & 0.0188 & 0.0050 & 0.0475 & 0.0000 & 2.3995 & 840 & 12.2 \\
\midrule
\multirow{3}{*}{\texttt{v02}} & 6 & 0.0133 & \textbf{0.0001} & 0.0133 & \textbf{0.0001} & 0.0012 & 0.0000 & 0.0589 & 792 & 4.1 \\
 & 14 & 0.0072 & 0.0006 & 0.0072 & 0.0006 & 0.0020 & 0.0000 & 0.0728 & 658 & 8.0 \\
 & 26 & 0.0045 & 0.0017 & 0.0045 & 0.0017 & 0.0032 & 0.0000 & 0.1203 & 828 & 15.0 \\
\midrule
\multirow{3}{*}{\texttt{v03}} & 6 & 0.1472 & 0.0552 & 0.1473 & 0.0554 & 0.0300 & 0.0000 & 6.0488 & 843 & 3.0 \\
 & 14 & 0.0608 & 0.0287 & 0.0610 & 0.0287 & 0.0404 & 0.0000 & 0.9403 & 840 & 6.7 \\
 & 26 & 0.0247 & 0.0070 & 0.0248 & 0.0070 & 0.0578 & 0.0000 & 2.8830 & 823 & 12.3 \\
\midrule
\multirow{3}{*}{\texttt{v04}} & 6 & 0.0062 & \textbf{0.0001} & 0.0062 & \textbf{0.0001} & \textbf{0.0008} & 0.0000 & 0.0604 & 648 & 4.0 \\
 & 14 & 0.0056 & 0.0007 & 0.0056 & 0.0007 & 0.0014 & 0.0000 & 0.0632 & 641 & 8.0 \\
 & 26 & \textbf{0.0032} & 0.0018 & \textbf{0.0032} & 0.0018 & 0.0018 & 0.0000 & 0.0526 & 756 & 15.1 \\
\midrule
\multirow{3}{*}{\texttt{v05}} & 6 & 0.2456 & 0.0712 & 0.2468 & 0.0712 & 0.0831 & 0.0000 & 45.1723 & 839 & 3.2 \\
 & 14 & 0.0507 & 0.0279 & 0.0508 & 0.0281 & 0.0530 & 0.0000 & 0.2993 & 812 & 6.9 \\
 & 26 & 0.0299 & 0.0128 & 0.0302 & 0.0128 & 0.0526 & 0.0000 & 1.3470 & \textbf{848} & 12.7 \\
\bottomrule
\end{tabular*}

\vspace{1.5ex}
\begin{minipage}{\textwidth}
Nota: Se predice un plano por objeto y la normal se elige con $S_{dat}$. Imputación y cobertura según la Sección~\ref{sec:sm_cobertura}.
\end{minipage}
\end{table}






\newpage
\clearpage
\subsection{Flujo A: Variantes de los prompts base}

\begin{table}[H]
\centering
\caption{Métricas de las variantes \texttt{\_1} del Flujo A en simetría planar, resolución 224px, iluminación \textit{flat}, envío multiimagen.}
\label{tab:nov_plane_variantes_1}
\footnotesize
\setlength{\tabcolsep}{3pt}
\renewcommand{\arraystretch}{0.95}

\begin{tabular*}{\textwidth}{@{\extracolsep{\fill}}lrrrrrrrr@{}}
\toprule
\multirow{2}{*}{Prompt} & \multirow{2}{*}{Vistas} & \multirow{2}{*}{AUC$_{45^\circ}$} & \multicolumn{3}{c}{Error angular ($^\circ$)} & \multicolumn{3}{c}{Precisión} \\
\cmidrule(lr){4-6} \cmidrule(lr){7-9}
 &  &  & Media & Mediana & Desv. & @5$^\circ$ & @10$^\circ$ & @15$^\circ$ \\
\midrule
\multirow{3}{*}{\texttt{v00\_1}} & 6 & 0.740 & 20.5 & 3.0 & 34.4 & 0.731 & \textbf{0.785} & 0.791 \\
 & 14 & 0.554 & 36.4 & 4.9 & 40.6 & 0.511 & 0.600 & 0.612 \\
 & 26 & \textbf{0.794} & \textbf{12.8} & \textbf{2.7} & 24.1 & 0.738 & 0.779 & \textbf{0.817} \\
\midrule
\multirow{3}{*}{\texttt{v01\_1}} & 6 & 0.153 & 71.3 & 89.2 & 31.7 & 0.133 & 0.149 & 0.155 \\
 & 14 & 0.028 & 79.2 & 84.5 & 15.8 & 0.008 & 0.022 & 0.029 \\
 & 26 & 0.017 & 87.3 & 89.6 & 11.9 & 0.017 & 0.018 & 0.018 \\
\midrule
\multirow{3}{*}{\texttt{v02\_1}} & 6 & 0.709 & 23.2 & 3.1 & 36.2 & 0.705 & 0.751 & 0.754 \\
 & 14 & 0.512 & 40.4 & 5.8 & 41.5 & 0.444 & 0.547 & 0.566 \\
 & 26 & 0.727 & 18.8 & \textbf{2.7} & 30.7 & 0.680 & 0.712 & 0.744 \\
\midrule
\multirow{3}{*}{\texttt{v03\_1}} & 6 & 0.159 & 70.0 & 87.9 & 32.2 & 0.148 & 0.160 & 0.165 \\
 & 14 & 0.037 & 78.4 & 85.0 & 18.0 & 0.019 & 0.031 & 0.040 \\
 & 26 & 0.020 & 85.4 & 89.5 & 13.9 & 0.013 & 0.018 & 0.020 \\
\midrule
\multirow{3}{*}{\texttt{v04\_1}} & 6 & 0.717 & 22.9 & 3.0 & 36.3 & \textbf{0.739} & 0.761 & 0.763 \\
 & 14 & 0.544 & 37.7 & 4.9 & 41.2 & 0.522 & 0.592 & 0.607 \\
 & 26 & 0.671 & 24.9 & \textbf{2.7} & 36.1 & 0.668 & 0.679 & 0.689 \\
\midrule
\multirow{3}{*}{\texttt{v05\_1}} & 6 & 0.631 & 28.8 & 4.0 & 37.9 & 0.584 & 0.663 & 0.673 \\
 & 14 & 0.540 & 36.0 & 6.2 & 39.4 & 0.413 & 0.556 & 0.587 \\
 & 26 & 0.603 & 25.3 & 9.1 & 31.1 & 0.448 & 0.519 & 0.601 \\
\bottomrule
\end{tabular*}

\vspace{2ex}

\begin{tabular*}{\textwidth}{@{\extracolsep{\fill}}lrrrrrrrrrr@{}}
\toprule
\multirow{2}{*}{Prompt} & \multirow{2}{*}{Vistas} & \multicolumn{2}{c}{Err.\ trasl.} & \multicolumn{2}{c}{Err.\ trasl.\ norm.} & \multicolumn{3}{c}{SDE$_{\mathrm{ref}}$} & \multicolumn{2}{c}{Cobertura} \\
\cmidrule(lr){3-4} \cmidrule(lr){5-6} \cmidrule(lr){7-9} \cmidrule(lr){10-11}
 &  & Media & Mediana & Media & Mediana & Media & Mín. & Máx. & Objetos & Puntos \\
\midrule
\multirow{3}{*}{\texttt{v00\_1}} & 6 & 0.0116 & \textbf{0.0001} & 0.0117 & \textbf{0.0001} & 0.0012 & 0.0000 & 0.0861 & 737 & 4.0 \\
 & 14 & 0.0084 & 0.0006 & 0.0085 & 0.0006 & 0.0021 & 0.0000 & 0.0771 & 582 & 8.0 \\
 & 26 & 0.0054 & 0.0018 & 0.0055 & 0.0018 & 0.0034 & 0.0000 & 0.3344 & 837 & 15.0 \\
\midrule
\multirow{3}{*}{\texttt{v01\_1}} & 6 & 0.1350 & 0.0547 & 0.1366 & 0.0547 & 0.0150 & 0.0000 & 0.3084 & \textbf{847} & 3.1 \\
 & 14 & 0.0459 & 0.0224 & 0.0460 & 0.0227 & 0.0323 & 0.0000 & 0.2910 & \textbf{847} & 6.6 \\
 & 26 & 0.0185 & 0.0059 & 0.0185 & 0.0059 & 0.0400 & 0.0000 & 0.9435 & 844 & 12.2 \\
\midrule
\multirow{3}{*}{\texttt{v02\_1}} & 6 & 0.0119 & \textbf{0.0001} & 0.0120 & \textbf{0.0001} & 0.0013 & 0.0000 & 0.0777 & 717 & 4.0 \\
 & 14 & 0.0108 & 0.0007 & 0.0108 & 0.0007 & 0.0027 & 0.0000 & 0.1518 & 544 & 8.0 \\
 & 26 & 0.0051 & 0.0018 & 0.0052 & 0.0018 & 0.0032 & 0.0000 & 0.1134 & 779 & 15.0 \\
\midrule
\multirow{3}{*}{\texttt{v03\_1}} & 6 & 0.1415 & 0.0500 & 0.1417 & 0.0501 & 0.0239 & 0.0000 & 0.5415 & 838 & 3.0 \\
 & 14 & 0.0588 & 0.0269 & 0.0589 & 0.0269 & 0.0381 & 0.0000 & 0.3278 & 841 & 6.6 \\
 & 26 & 0.0219 & 0.0081 & 0.0219 & 0.0081 & 0.0444 & 0.0000 & 2.1872 & 826 & 12.3 \\
\midrule
\multirow{3}{*}{\texttt{v04\_1}} & 6 & 0.0059 & \textbf{0.0001} & 0.0060 & \textbf{0.0001} & \textbf{0.0009} & 0.0000 & 0.0685 & 707 & 4.1 \\
 & 14 & 0.0049 & 0.0006 & 0.0049 & 0.0006 & 0.0015 & 0.0000 & 0.0662 & 553 & 8.0 \\
 & 26 & \textbf{0.0032} & 0.0018 & \textbf{0.0032} & 0.0018 & 0.0021 & 0.0000 & 0.1109 & 687 & 15.1 \\
\midrule
\multirow{3}{*}{\texttt{v05\_1}} & 6 & 0.0591 & 0.0019 & 0.0592 & 0.0019 & 0.0079 & 0.0000 & 0.3240 & 796 & 3.9 \\
 & 14 & 0.0171 & 0.0011 & 0.0171 & 0.0011 & 0.0111 & 0.0000 & 0.2697 & 684 & 7.8 \\
 & 26 & 0.0147 & 0.0019 & 0.0148 & 0.0019 & 0.0168 & 0.0000 & 1.2787 & 827 & 14.8 \\
\bottomrule
\end{tabular*}

\vspace{1.5ex}
\begin{minipage}{\textwidth}
Nota: Se predice un plano por objeto y la normal se elige con $S_{dat}$. Cada prompt es una variante del prompt base con la que se busca mejorar sus resultados. Imputación y cobertura según la Sección~\ref{sec:sm_cobertura}.
\end{minipage}
\end{table}



\newpage
\clearpage
\subsection{Flujo A: Prompts adicionales}

\begin{table}[H]
\centering
\caption{Métricas de los prompts adicionales del Flujo A en simetría planar, resolución 224px, iluminación \textit{flat}, envío multiimagen.}
\label{tab:nov_plane_adicionales}
\footnotesize
\setlength{\tabcolsep}{3pt}
\renewcommand{\arraystretch}{0.95}

\begin{tabular*}{\textwidth}{@{\extracolsep{\fill}}lrrrrrrrr@{}}
\toprule
\multirow{2}{*}{Prompt} & \multirow{2}{*}{Vistas} & \multirow{2}{*}{AUC$_{45^\circ}$} & \multicolumn{3}{c}{Error angular ($^\circ$)} & \multicolumn{3}{c}{Precisión} \\
\cmidrule(lr){4-6} \cmidrule(lr){7-9}
 &  &  & Media & Mediana & Desv. & @5$^\circ$ & @10$^\circ$ & @15$^\circ$ \\
\midrule
\multirow{3}{*}{\texttt{v06}} & 6 & 0.726 & 21.6 & 3.1 & 35.1 & 0.714 & \textbf{0.767} & 0.772 \\
 & 14 & 0.510 & 40.9 & 6.0 & 41.8 & 0.436 & 0.540 & 0.564 \\
 & 26 & \textbf{0.762} & \textbf{14.9} & \textbf{2.6} & 26.1 & 0.721 & 0.753 & \textbf{0.779} \\
\midrule
\multirow{3}{*}{\texttt{v07}} & 6 & 0.658 & 27.9 & 3.1 & 38.8 & 0.665 & 0.695 & 0.699 \\
 & 14 & 0.490 & 43.0 & 5.1 & 42.2 & 0.480 & 0.541 & 0.545 \\
 & 26 & 0.710 & 22.5 & 2.7 & 35.5 & \textbf{0.729} & 0.734 & 0.738 \\
\midrule
\multirow{3}{*}{\texttt{v08}} & 6 & 0.199 & 65.1 & 85.5 & 34.1 & 0.167 & 0.188 & 0.195 \\
 & 14 & 0.097 & 71.4 & 84.3 & 25.7 & 0.045 & 0.084 & 0.095 \\
 & 26 & 0.099 & 75.9 & 88.3 & 26.5 & 0.057 & 0.091 & 0.106 \\
\bottomrule
\end{tabular*}

\vspace{2ex}

\begin{tabular*}{\textwidth}{@{\extracolsep{\fill}}lrrrrrrrrrr@{}}
\toprule
\multirow{2}{*}{Prompt} & \multirow{2}{*}{Vistas} & \multicolumn{2}{c}{Err.\ trasl.} & \multicolumn{2}{c}{Err.\ trasl.\ norm.} & \multicolumn{3}{c}{SDE$_{\mathrm{ref}}$} & \multicolumn{2}{c}{Cobertura} \\
\cmidrule(lr){3-4} \cmidrule(lr){5-6} \cmidrule(lr){7-9} \cmidrule(lr){10-11}
 &  & Media & Mediana & Media & Mediana & Media & Mín. & Máx. & Objetos & Puntos \\
\midrule
\multirow{3}{*}{\texttt{v06}} & 6 & 0.0108 & \textbf{0.0001} & 0.0111 & \textbf{0.0001} & \textbf{0.0011} & 0.0000 & 0.0629 & 730 & 4.0 \\
 & 14 & 0.0091 & 0.0007 & 0.0092 & 0.0007 & 0.0020 & 0.0000 & 0.1006 & 531 & 8.0 \\
 & 26 & 0.0056 & 0.0018 & 0.0056 & 0.0018 & 0.0030 & 0.0000 & 0.0980 & 822 & 15.0 \\
\midrule
\multirow{3}{*}{\texttt{v07}} & 6 & 0.0053 & \textbf{0.0001} & 0.0053 & \textbf{0.0001} & 0.0012 & 0.0000 & 0.0829 & 652 & 4.1 \\
 & 14 & 0.0049 & 0.0006 & 0.0049 & 0.0006 & 0.0012 & 0.0000 & 0.0568 & 501 & 8.0 \\
 & 26 & \textbf{0.0038} & 0.0018 & \textbf{0.0038} & 0.0018 & 0.0017 & 0.0000 & 0.1100 & 709 & 15.0 \\
\midrule
\multirow{3}{*}{\texttt{v08}} & 6 & 0.2293 & 0.0661 & 0.2296 & 0.0661 & 0.0413 & 0.0000 & 4.9261 & 841 & 3.2 \\
 & 14 & 0.0736 & 0.0316 & 0.0736 & 0.0316 & 0.0515 & 0.0000 & 0.4469 & \textbf{844} & 7.0 \\
 & 26 & 0.0310 & 0.0138 & 0.0310 & 0.0138 & 0.0555 & 0.0000 & 2.2473 & \textbf{844} & 12.8 \\
\bottomrule
\end{tabular*}

\vspace{1.5ex}
\begin{minipage}{\textwidth}
Nota: Se predice un plano por objeto y la normal se elige con $S_{dat}$. Imputación y cobertura según la Sección~\ref{sec:sm_cobertura}.
\end{minipage}
\end{table}





\newpage
\subsection{Flujo A: Top}
\label{subsec:plane_top_224}

Se aplica el mismo procedimiento descrito para la simetría axial (Sección~\ref{subsec:axis_top_224}). El orden de prioridad, con la brecha entre paréntesis, es: (1) \textbf{$F1_{\text{ref}}$ voraz} ($0.05$); (2) \textbf{AUC} ($0.05$), que compensa que $F1_{\text{ref}}$ se calcule solo sobre los objetos con predicción, ya que imputa $90°$ a los demás; (3) \textbf{P@$10°$} ($0.05$); (4) \textbf{cobertura} (17 objetos, 2\% de 850); (5) \textbf{SDE} ($0.005$); y (6) \textbf{error angular medio}, sin brecha.

La variante húngara de $F1_{\text{ref}}$ no se usa para desempatar porque, con un solo plano predicho por objeto, ordena las configuraciones casi igual que la voraz. El error angular mediano tampoco se incluye: la distribución del error es bimodal (aciertos cercanos a $0°$ y planos perpendiculares cercanos a $90°$), de modo que la mediana salta entre ambos modos, y esa información ya está contenida en P@$10°$. La Tabla~\ref{tab:plane_top} muestra las tres mejores configuraciones.



Como en simetría axial, la selección compara los promedios de cada \textit{prompt} sobre 6, 14 y 26 vistas (Tabla~\ref{tab:plane_top_prompts}), con una diferencia: el AUC pasa a ser el primer criterio y $F1_{\text{ref}}$ el segundo. Al repetir un \textit{prompt} con todas las vistas, su cobertura en cada una debe pesar, y $F1_{\text{ref}}$ no la considera. Por ejemplo, \texttt{v07} lidera el ranking con 26 vistas, pero con 14 vistas predice solo 501 objetos y su AUC baja a 0.490.



\texttt{v00\_1}, \texttt{v04} y \texttt{v02} empatan en AUC. \texttt{v02} queda fuera del desempate del primer puesto por $F1_{\text{ref}}$ (0.600 frente a 0.653, apenas sobre la brecha), y \texttt{v00\_1} supera a \texttt{v04} por cobertura. Se seleccionan estos tres \textit{prompts}, sin que su orden indique una superioridad clara: \texttt{v02} obtiene el mejor AUC, P@$10°$, cobertura y error medio. Además, usan estrategias distintas: traza del plano, puntos medios y costura, respectivamente.


% Generado por utils/generar_tablas_seleccion.py. No editar a mano.
\begin{table}[H]
\centering
\caption{Promedio por \textit{prompt} sobre 6, 14 y 26 vistas, en las mismas condiciones.}
\label{tab:plane_top_prompts}
\footnotesize
\setlength{\tabcolsep}{3pt}
\renewcommand{\arraystretch}{0.95}

\begin{tabular*}{\textwidth}{@{\extracolsep{\fill}}rlcrrrrrr@{}}
\toprule
\# & Prompt & Crit. & AUC$_{45^\circ}$ & \begin{tabular}[b]{@{}c@{}}$F1_{\mathrm{ref}}$ \\ voraz\end{tabular} & P@10$^\circ$ & Objetos & SDE$_{\mathrm{ref}}$ & \begin{tabular}[b]{@{}c@{}}Err.\ medio \\ ($^\circ$)\end{tabular} \\
\midrule
1 & \texttt{v00\_1} & 4 & 0.696 & 0.615 & 0.721 & 719 & 0.0022 & 23.2 \\
2 & \texttt{v04} & 2 & 0.678 & 0.653 & 0.702 & 682 & \textbf{0.0014} & 25.2 \\
3 & \texttt{v02} & 1 & \textbf{0.727} & 0.600 & \textbf{0.751} & \textbf{759} & 0.0022 & \textbf{19.9} \\
4 & \texttt{v04\_1} & 4 & 0.644 & \textbf{0.658} & 0.677 & 649 & 0.0015 & 28.5 \\
\bottomrule
\end{tabular*}

\vspace{0.5ex}
\begin{minipage}{\textwidth}
Nota: Crit.: criterio que resolvió la posición, según el orden de las columnas.
\end{minipage}
\end{table}

\texttt{v00\_1}, \texttt{v04} y \texttt{v02} empatan en AUC. \texttt{v02} queda fuera del desempate del primer puesto por $F1_{\text{ref}}$ (0.600 frente a 0.653, apenas sobre la brecha), y \texttt{v00\_1} supera a \texttt{v04} por cobertura. Se seleccionan estos tres \textit{prompts}, sin que su orden indique una superioridad clara: \texttt{v02} obtiene el mejor AUC, P@$10°$, cobertura y error medio. Además, usan estrategias distintas: traza del plano, puntos medios y costura, respectivamente.


\newpage
\subsection{Flujo A: Experimento de resolución}

Se utilizan los prompts \texttt{v00\_1}, \texttt{v04} y \texttt{v02} en resoluciones 224px, 448px y 896px, con iluminación \textit{flat} y envío multiimagen.

\begin{table}[H]
\centering
\caption{Métricas de los prompts seleccionados en simetría planar, resolución 448px, iluminación \textit{flat}, envío multiimagen.}
\label{tab:nov_plane_448}
\footnotesize
\setlength{\tabcolsep}{3pt}
\renewcommand{\arraystretch}{0.95}

\begin{tabular*}{\textwidth}{@{\extracolsep{\fill}}lrrrrrrrr@{}}
\toprule
\multirow{2}{*}{Prompt} & \multirow{2}{*}{Vistas} & \multirow{2}{*}{AUC$_{45^\circ}$} & \multicolumn{3}{c}{Error angular ($^\circ$)} & \multicolumn{3}{c}{Precisión} \\
\cmidrule(lr){4-6} \cmidrule(lr){7-9}
 &  &  & Media & Mediana & Desv. & @5$^\circ$ & @10$^\circ$ & @15$^\circ$ \\
\midrule
\multirow{3}{*}{\texttt{v00\_1}} & 6 & 0.640 & 29.9 & 3.2 & 39.9 & 0.631 & 0.672 & 0.682 \\
 & 14 & 0.523 & 39.7 & 5.0 & 41.6 & 0.494 & 0.566 & 0.581 \\
 & 26 & 0.809 & 13.1 & 2.7 & 26.6 & 0.833 & 0.842 & 0.845 \\
\midrule
\multirow{3}{*}{\texttt{v04}} & 6 & 0.484 & 44.8 & 6.4 & 43.4 & 0.491 & 0.514 & 0.516 \\
 & 14 & 0.442 & 47.5 & 64.8 & 42.3 & 0.428 & 0.483 & 0.491 \\
 & 26 & 0.792 & 12.9 & 2.7 & 24.4 & 0.795 & 0.802 & 0.805 \\
\midrule
\multirow{3}{*}{\texttt{v02}} & 6 & 0.742 & 20.2 & 3.1 & 34.1 & 0.722 & 0.783 & 0.791 \\
 & 14 & 0.536 & 38.1 & 5.1 & 41.1 & 0.481 & 0.579 & 0.591 \\
 & 26 & \textbf{0.827} & \textbf{12.3} & \textbf{2.6} & 26.6 & \textbf{0.839} & \textbf{0.857} & \textbf{0.866} \\
\bottomrule
\end{tabular*}

\vspace{2ex}

\begin{tabular*}{\textwidth}{@{\extracolsep{\fill}}lrrrrrrrrrr@{}}
\toprule
\multirow{2}{*}{Prompt} & \multirow{2}{*}{Vistas} & \multicolumn{2}{c}{Err.\ trasl.} & \multicolumn{2}{c}{Err.\ trasl.\ norm.} & \multicolumn{3}{c}{SDE$_{\mathrm{ref}}$} & \multicolumn{2}{c}{Cobertura} \\
\cmidrule(lr){3-4} \cmidrule(lr){5-6} \cmidrule(lr){7-9} \cmidrule(lr){10-11}
 &  & Media & Mediana & Media & Mediana & Media & Mín. & Máx. & Objetos & Puntos \\
\midrule
\multirow{3}{*}{\texttt{v00\_1}} & 6 & 0.0108 & \textbf{0.0001} & 0.0110 & \textbf{0.0001} & 0.0019 & 0.0000 & 0.1326 & 642 & 4.1 \\
 & 14 & 0.0076 & 0.0007 & 0.0077 & 0.0007 & 0.0024 & 0.0000 & 0.1685 & 543 & 8.0 \\
 & 26 & 0.0038 & 0.0019 & 0.0038 & 0.0019 & 0.0017 & 0.0000 & 0.0706 & 820 & 15.0 \\
\midrule
\multirow{3}{*}{\texttt{v04}} & 6 & 0.0057 & \textbf{0.0001} & 0.0057 & \textbf{0.0001} & \textbf{0.0011} & 0.0000 & 0.0682 & 464 & 4.1 \\
 & 14 & 0.0065 & 0.0007 & 0.0065 & 0.0007 & 0.0017 & 0.0000 & 0.0787 & 460 & 8.0 \\
 & 26 & \textbf{0.0031} & 0.0019 & \textbf{0.0032} & 0.0019 & 0.0021 & 0.0000 & 0.0758 & \textbf{827} & 15.1 \\
\midrule
\multirow{3}{*}{\texttt{v02}} & 6 & 0.0122 & \textbf{0.0001} & 0.0123 & \textbf{0.0001} & 0.0495 & 0.0000 & 35.7678 & 752 & 4.0 \\
 & 14 & 0.0073 & 0.0007 & 0.0073 & 0.0007 & 0.0018 & 0.0000 & 0.0905 & 563 & 8.0 \\
 & 26 & 0.0034 & 0.0018 & 0.0034 & 0.0018 & 0.0014 & 0.0000 & 0.0694 & 815 & 15.0 \\
\bottomrule
\end{tabular*}

\vspace{1.5ex}
\begin{minipage}{\textwidth}
Nota: Se predice un plano por objeto y la normal se elige con $S_{dat}$. Imputación y cobertura según la Sección~\ref{sec:sm_cobertura}.
\end{minipage}
\end{table}

\newpage

\clearpage

\begin{table}[htbp]
\centering
\caption{Noviembre: F1-scores de los prompts seleccionados en simetría planar, resolución 448px, iluminación \textit{flat}, envío multiimagen (\textit{multi}).}
\label{tab:nov_plane_448_f1}
\begin{lrbox}{\tablebox}
\begin{adjustbox}{max width=\textwidth}
\begin{tabular}{lrrrrrrrrrrr}
\toprule
\multirow{2}{*}{Prompt} & \multirow{2}{*}{Vistas} & \multicolumn{5}{c}{$F1_{\mathrm{ref}}$ voraz} & \multicolumn{5}{c}{$F1_{\mathrm{ref}}$ húngaro} \\
\cmidrule(lr){3-7} \cmidrule(lr){8-12}
 &  & Prom. & $\tau{=}0.05$ & $\tau{=}0.10$ & $\tau{=}0.15$ & $\tau{=}0.20$ & Prom. & $\tau{=}0.05$ & $\tau{=}0.10$ & $\tau{=}0.15$ & $\tau{=}0.20$ \\
\midrule
\multirow{3}{*}{\texttt{v00\_1}} & 6 & 0.651 & 0.367 & 0.723 & 0.752 & 0.761 & 0.705 & 0.397 & 0.783 & 0.814 & 0.824 \\
 & 14 & 0.541 & 0.032 & 0.679 & 0.722 & 0.731 & 0.594 & 0.035 & 0.746 & 0.792 & 0.802 \\
 & 26 & 0.716 & 0.671 & 0.725 & 0.734 & 0.735 & 0.778 & 0.729 & 0.788 & 0.798 & 0.799 \\
\midrule
\multirow{3}{*}{\texttt{v04}} & 6 & 0.702 & 0.411 & \textbf{0.781} & \textbf{0.804} & \textbf{0.811} & 0.755 & 0.442 & \textbf{0.841} & \textbf{0.865} & \textbf{0.873} \\
 & 14 & 0.568 & 0.058 & 0.706 & 0.745 & 0.764 & 0.613 & 0.062 & 0.761 & 0.804 & 0.824 \\
 & 26 & 0.674 & 0.633 & 0.683 & 0.690 & 0.690 & 0.734 & 0.689 & 0.745 & 0.751 & 0.751 \\
\midrule
\multirow{3}{*}{\texttt{v02}} & 6 & 0.634 & 0.356 & 0.693 & 0.740 & 0.748 & 0.689 & 0.387 & 0.753 & 0.804 & 0.812 \\
 & 14 & 0.540 & 0.045 & 0.661 & 0.723 & 0.732 & 0.588 & 0.049 & 0.719 & 0.787 & 0.797 \\
 & 26 & \textbf{0.730} & \textbf{0.688} & 0.734 & 0.745 & 0.753 & \textbf{0.794} & \textbf{0.749} & 0.798 & 0.810 & 0.819 \\
\bottomrule
\end{tabular}
\end{adjustbox}
\end{lrbox}
\usebox{\tablebox}

\vspace{1.5ex}
\begin{minipage}{\wd\tablebox}
\footnotesize
Nota: La normal se elige con $S_{dat}$ y se predice un plano por objeto. Un plano predicho es acierto si su distancia a un plano de referencia es menor que $\tau$. Prom.: promedio de los cuatro umbrales. Voraz y húngaro difieren en cómo se emparejan los planos predichos con los de referencia. Como se predice un solo plano por objeto, en objetos con varias simetrías no se pueden detectar todas. Solo se usan los objetos con predicción. Negrita: mejor valor de la columna.
\end{minipage}
\end{table}



\newpage


























\newpage
\subsection{Flujo A: Experimento de iluminación}

Se repite el experimento con \texttt{v05\_1}, el \textit{prompt} planar con mejor
resultado individual a 224px (AUC $=0.494$ con 6 vistas), bajo las mismas tres
condiciones y prediciendo un plano por objeto, igual que en la corrida
\textit{flat}.

Nota: Ojo con la redacción y las citas. En \texttt{ImagesGenerator/export\_fibonacci\_views.py}, la ``iluminación'' se implementa cambiando el color gris de la malla (TexturesVertex), con la luz direccional fija. No cambia la intensidad de la luz. Por eso escribí ``tono de gris uniforme de la malla'' y no ``intensidad de la luz''. \texttt{lectura\_iluminacion.md} advierte que 3D Common Corruptions solo aplica si se modifica la luz de la escena; con esta implementación no correspondería citarlo. VLM-RobustBench, que evalúa Molmo ante perturbaciones fotométricas, sí sirve, pero tiene la venue pendiente de verificar.






\newpage
\subsection{Flujo B y C}

Se aplica el mismo diseño con \texttt{v05\_1}. En el Flujo C, la pista describe por
dónde pasa la traza del plano sobre la superficie, y la corrida usa
\texttt{v04\_1} como respaldo.







\newpage

\subsection{Flujo A: Envío individual vs Multiimagen}

Al aumentar el número de vistas, el error del modo multiimagen crece. Este
experimento prueba si la causa es la cantidad de imágenes en una misma llamada. En
el modo multiimagen, las $n_v$ vistas se envían juntas (con 26 vistas, unos 11\,000
tokens de entrada). En el modo individual, cada vista se envía en una llamada
aparte, con el \textit{prompt} de imagen única, y los puntos se reúnen antes de la
triangulación. Así, la predicción de cada vista no depende de cuántas otras vistas
haya: si el error sigue creciendo con $n_v$, la causa está en las vistas o en el
estimador, no en la saturación del modelo. Las dos variantes difieren también en el
texto del \textit{prompt}, porque la versión de imagen única no pide coherencia
entre vistas.




\newpage
\subsection{Baseline multiplano}







\newpage
\section{Diagnósticos}


\subsection{Colapso al centro de la imagen}
\label{sec:sm_diag_centro}

El diagnóstico \texttt{diagnose\_center\_bias.py} mide, por prompt y $n_v$, la
fracción de puntos con $|x-500|<30$ en unidades de Molmo2, separando las vistas
cenitales (elevación $>70°$ en valor absoluto, donde $x\approx 500$ es la respuesta
correcta) de las demás. La Tabla~\ref{tab:sm_centerbias} muestra que el colapso es
generalizado: 40--99\% en la mayoría de los prompts, con la excepción de los
bilaterales, \texttt{plane\_v05} y los \texttt{v08} (4--25\%).

\begin{table}[htbp]
\centering
\caption{Correlación entre el porcentaje de centrado y el desempeño (57 filas
prompt $\times$ $n_v$ por tipo).}
\label{tab:sm_centro_corr}
\resizebox{\textwidth}{!}{%
\begin{tabular}{llrrl}
\toprule
Tipo & Métrica & Pearson & Spearman & Lectura \\
\midrule
Eje & error angular medio & $+0.13$ & $-0.10$ & sin relación \\
Eje & AUC & $-0.12$ & $0.00$ & sin relación \\
Eje & traslación normalizada & $\mathbf{+0.72}$ & $\mathbf{+0.79}$ & más centrado, ancla peor condicionada \\
Plano & $R_{adj}$ & $\mathbf{+0.75}$ & $\mathbf{+0.72}$ & más centrado, mejor \\
Plano & $F1_{adj}$ & $\mathbf{+0.81}$ & $\mathbf{+0.74}$ & más centrado, mejor \\
Plano & $F1_{\text{ref}}$ & $\mathbf{+0.67}$ & $\mathbf{+0.67}$ & más centrado, mejor \\
\bottomrule
\end{tabular}}
\end{table}

El mismo comportamiento de Molmo2 es inútil en eje y ``exitoso'' en plano
(Tabla~\ref{tab:sm_centro_corr}). La sección siguiente explica por qué.

\subsection{Orientación de la referencia y colapso a los ejes del mundo}
\label{sec:sm_orientacion}

Se analizó la orientación de la referencia en los 30+30 objetos disponibles
localmente, y se re-ejecutó el código de producción sobre sus puntos de Molmo2.

\begin{table}[htbp]
\centering
\caption{Orientación de los vectores de referencia (muestra local de 30 objetos por tipo).}
\label{tab:sm_gt_orient}
\begin{tabular}{p{5.2cm}p{5.4cm}p{3.4cm}}
\toprule
 & Eje & Plano \\
\midrule
A menos de $5°$ de un eje del mundo & 10\% & \textbf{100\%} \\
A menos de $15°$ de un eje del mundo & 27\% & \textbf{100\%} \\
Distribución & aleatoria (media de $|\hat{\mathbf{d}}\cdot\hat{Y}|$: 0.45; uniforme: 0.50) & 23 normales $X$, 8 normales $Z$ \\
Origen de la referencia & $(0,0,0)$ en todos & $(0,0,0)$ en todos \\
\bottomrule
\end{tabular}
\end{table}

Los objetos axiales fueron rotados al azar; los planares conservan la pose canónica
de ShapeNet (Tabla~\ref{tab:sm_gt_orient}). Por otro lado, el renderizado usa
\texttt{look\_at\_view\_transform}: todas las cámaras miran al origen con ``arriba''
$=+Y$, y para un $n_v$ dado el conjunto de cámaras es el mismo para todos los objetos.

\paragraph{Eje.} Re-ejecutando la triangulación sobre \texttt{axis\_v06} en la
muestra, 12 de 29 ejes predichos (con 26 vistas) son exactamente $+Y$ (a
$0.0°$--$1.2°$), y el 48\% de las predicciones queda a menos de $5°$ de algún eje del
mundo (frente al 10\% de las referencias). Los casos ``buenos'' del sandbox son
objetos cuyo eje de referencia casualmente es vertical: con $|\hat{\mathbf{d}}_{gt}\cdot\hat{Y}|\approx 1$
el error es $2$--$6°$; con la referencia horizontal, $87$--$90°$. Predecir siempre la
vertical produce, frente a referencias orientadas al azar, exactamente el error del
azar.

\paragraph{Plano.} Cuando Molmo2 marca la columna central, las rectas trianguladas
son verticales y las normales candidatas quedan en el plano horizontal, alineadas
con $X$ o $Z$ del mundo. Con referencias canónicas mayoritariamente $X$, eso acierta
sin que Molmo2 haya identificado la simetría del objeto. En la muestra, el
pipeline de plano (\texttt{plane\_v04\_1}, 14 vistas) predice 1.67 planos por
objeto y 27 de sus 50 normales están alineadas con $X$.

\subsection{Predictores triviales}
\label{sec:sm_triviales}

\begin{table}[htbp]
\centering
\caption{Predictores que ignoran la imagen, evaluados con las funciones de
\texttt{evaluate.py} sobre la muestra local de 30+30 objetos.}
\label{tab:sm_triviales}
\begin{tabular}{lccccc}
\toprule
Predictor & \multicolumn{2}{c}{Eje} & \multicolumn{3}{c}{Plano} \\
\cmidrule(lr){2-3}\cmidrule(lr){4-6}
 & $\bar{\varepsilon}$ & AUC & $R$ / $P$ & $F1_{\text{ref}}$ & $F1^{H}_{\text{ref}}$ \\
\midrule
Azar uniforme (analítico) & 57.30 & 0.100 & -- & -- & -- \\
Siempre $X$ & 55.10 & 0.128 & 0.767 / 0.767 & \textbf{0.742} & 0.754 \\
Siempre $Y$ & 58.53 & 0.163 & 0 / 0 & 0 & 0 \\
Siempre $Z$ & 60.07 & 0.094 & 0.217 / 0.233 & 0.258 & 0.230 \\
Siempre $\{X,Z\}$ & -- & -- & 0.983 / 0.500 & 0.886 & 0.659 \\
Siempre $\{X,Y,Z\}$ & -- & -- & 0.983 / 0.333 & 0.795 & 0.496 \\
\midrule
Pipeline, \texttt{plane\_v04\_1}, 14 vistas & -- & -- & 0.667 / 0.483 & 0.374 & -- \\
\bottomrule
\end{tabular}
\end{table}

En la muestra, el pipeline de plano queda por debajo de ``siempre $X$''
(Tabla~\ref{tab:sm_triviales}). Esta comparación debe repetirse sobre los 850
objetos antes de afirmar cualquier resultado de plano; el script correspondiente ya
existe y corre en segundos sin GPU.



\newpage
\section{Analisis multiplano}

Para simetria planar (\texttt{max\_planes})

\section{Analisis Politicas de descarte de vistas}

Para simetria axial y planar

\section{Diagnostico de sesgo al centro x=500}

Para simetria axial y planar

\section{Condiciones de grilla}

Para simetria axial 

\section{Condiciones de grilla}

Para simetria planar (\texttt{dup\_angle\_thresh} y \texttt{sde\_gate})


\section{Comparación QWen vs Molmo2 vs MolmoPointing}

Esta sección evalúa si los resultados dependen del modelo. Se reemplaza Molmo2 por
[Qwen3-VL / modelo y versión] y [MolmoPointing: modelo y versión] sobre la mejor
configuración de Molmo2 de cada tipo de simetría ([prompt, resolución, vistas]),
manteniendo idénticos los renders, la triangulación y la evaluación. Además de las
métricas, se repite el diagnóstico de colapso al centro de la imagen
(Sección~\ref{sec:sm_diag_centro}), para determinar si el sesgo hacia
$x\approx500$ y el resultado en simetría axial son propios de Molmo2 o compartidos
por otros modelos con capacidad de \textit{pointing}.


\section{Resultados usando las categorias de los objetos (silla, pote, etc)}

\section{Seección especifica de vistas}

Best view-selection (para cada objeto) y comparación con resultados TOP del Flujo A

\newpage